% >>> PROPAGATION (Fable 5, Aug 29 2026): the Aetherian Nexus recorded as a payoff of null-set retention at the measure-zero
% passage (a reopening lifts the discount; the reopened blank weighted by the tree's own retained structure). See the ledger, R109.
% Audit and change history (289 lines) relocated verbatim to
% aethic_tree_changelog.md, August 12, 2026, per referee recommendation. In-text
% workflow tags de-jargonized the same day; refutation records remain by design.
\begin{filecontents}[overwrite]{alg_refs.bib}
% entries. Sibling-paper keys below; titles per current drafts.
% Sibling record codes are author-supplied from the site's paper directory (a
% dated communication logged in the session record); the gate should still
% click each URL before relying on it.
% MIGRATE AT MERGE: golan, tarski, and benci are new-to-the-corpus citations --- add to
% duplicates and may be removed.
@book{golan1999semirings, author={Golan, Jonathan S.}, title={Semirings and their Applications}, publisher={Kluwer}, year={1999}}
@article{tarski1955lattice, author={Tarski, Alfred}, title={A lattice-theoretical fixpoint theorem and its applications}, journal={Pacific Journal of Mathematics}, volume={5}, pages={285--309}, year={1955}}
@article{benci2013nonarchimedean, author={Benci, Vieri and Horsten, Leon and Wenmackers, Sylvia}, title={Non-Archimedean Probability}, journal={Milan Journal of Mathematics}, volume={81}, pages={121--151}, year={2013}, note = {Pages AI-confirmed (adversarial external pass, Aug 2026); human verification pending per convention}}
@misc{Aethus, howpublished={Manuscript}, author={Benander, Ajax}, title={Aethic Reasoning: A Comprehensive Solution to the Quantum Measurement Problem}, year={2024}, month = {11}, url={https://philpapers.org/rec/BENARA-5}, eprint={https://philsci-archive.pitt.edu/id/eprint/25713}, eprinttype={Alternative URL}, note={Manuscript; the complete development of the framework, successive deposits November 2024--2026}}
@misc{AethusBrief, author={Benander, Ajax}, title={Aethic Reasoning in Brief: From a Counterfactual Semantics to the Quantum Measurement Problem}, year={2026}, howpublished={Preprint}, url={https://philpapers.org/rec/BENARI-3}, eprint={https://philsci-archive.pitt.edu/id/eprint/30439}, eprinttype={Alternative URL}}

@misc{AethusUnion, author={Benander, Ajax}, title={Exclusivity as Achievement: Relocating Mutual Exclusion to the Interior, and the Composition Regime That Falls Out}, year={2026}, howpublished={Preprint}, url={https://philpapers.org/rec/BENEAA-9}}



@inproceedings{green2007provenance,
  author = {Green, Todd J. and Karvounarakis, Grigoris and Tannen, Val},
  title = {Provenance semirings},
  booktitle = {Proceedings of the 26th ACM Symposium on Principles of Database Systems},
  pages = {31--40}, year = {2007}
}
@book{droste2009handbook,
  editor = {Droste, Manfred and Kuich, Werner and Vogler, Heiko},
  title = {Handbook of Weighted Automata},
  publisher = {Springer}, address = {Berlin}, year = {2009}
}
@incollection{abramskyjung1994domain,
  author = {Abramsky, Samson and Jung, Achim},
  title = {Domain theory},
  booktitle = {Handbook of Logic in Computer Science},
  volume = {3}, pages = {1--168},
  publisher = {Oxford University Press}, address = {Oxford}, year = {1994}
}
@misc{AethusQuant,
  author       = {Benander, Ajax},
  title        = {Quantum Mechanics from Records: Measurement Structure on a Finite Path-Kernel Domain},
  year         = {2026},
  howpublished = {Preprint},
  url          = {https://philpapers.org/rec/BENTQB},
  note         = {The physics-face axiomatization; first posted July 17, 2026.}
}

@misc{AethusSemantics,
  author       = {Benander, Ajax},
  title        = {Counterfactual Evaluation by Perturbation of Information States: Deleting the World-Similarity Ordering, Not Refining It},
  year         = {2026},
  howpublished = {Preprint},
  url          = {https://philpapers.org/rec/BENCEB-2},
  note         = {The modal face; the Cautious Monotony countermodel and the
                  non-preferentiality theorem.}
}

@misc{NexicFoundation,
  author       = {Benander, Ajax},
  title        = {The Copernican Method Against the Copernican Principle: Anthropic Typicality over Partial Information States},
  year         = {2026},
  howpublished = {Preprint},
  url          = {https://philpapers.org/rec/BENTCM-7},
  note         = {The Nexic wing's foundation extraction; the Accordance Principle
                  and the statistical trilemma.}
}

@book{newton1687principia,
  author    = {Newton, Isaac},
  title     = {Philosophi\ae{} Naturalis Principia Mathematica},
  year      = {1687},
  address   = {London},
  publisher = {Jussu Societatis Regi\ae{} ac Typis Josephi Streater},
  note      = {The \emph{hypotheses non fingo} of the General Scholium was added in
               the second edition (Cambridge, 1713), ed.\ Roger Cotes; standard
               English text in the Cohen and Whitman translation,
               \emph{The Principia: Mathematical Principles of Natural Philosophy},
               University of California Press, 1999, p.~943.}
}

@misc{GoldenScience2022,
  author       = {Benander, Ajax},
  title        = {Golden Science},
  year         = {2022},
  month        = {November},
  howpublished = {Self-hosted manuscript},
  url          = {https://aethus.org/aethic-code/files/GoldenScience_Nov2022.pdf},
  note         = {Title page August 2022, the date of the first version of the source; this version created November 2022. Not deposited to a
                  preprint server; cited as a dated developmental record rather than
                  as a source of results.}
}
\end{filecontents}
\documentclass[11pt]{article}
\usepackage[margin=0.75in]{geometry}
\usepackage{amsmath, amssymb, amsthm}
\usepackage[hidelinks]{hyperref}
\usepackage{microtype}
\IfFileExists{lmodern.sty}{\usepackage{lmodern}}{\microtypesetup{expansion=false}}
\usepackage[style=numeric, sorting=none]{biblatex}
\addbibresource{refs.bib}
\addbibresource{alg_refs.bib}
\newtheorem{defin}{Definition}
\newtheorem{prop}{Proposition}
\newtheorem{lemma}{Lemma}
\newtheorem{theorem}{Theorem}
\newcommand{\Null}{\mathbf{1}}
\newcommand{\Hyp}{\mathbf{0}}

\title{The Weight Algebra of Aethae:\\ \large A Two-Layer Commutative Semiring for Superposition and Pruning}
\author{Ajax Benander\thanks{Email: \href{mailto:akb9408@rit.edu}{akb9408@rit.edu}, Website: \url{https://aethus.org}. This paper extracts, as a standalone algebraic development, the weighted-formulation section of the long companion manuscript \cite{Aethus}, keeping its constructions and their motivating remarks intact and stating them in the vocabulary of the companion's foundations. The fuller development gives the objects constructed here a proper name, \emph{Aethus} (plural \emph{Aethae}), and this paper uses it, stating first the premise the name is given to (Definition~\ref{def:aethus}).}}
\date{First drafted July 2026\\ {\small This version: September 2026 --- working draft}}

\usepackage{booktabs}
\begin{document}
\maketitle

\begin{abstract}
\noindent We construct the algebra formed by the partial information states of the companion framework: deductively closed partial assignments of values to attributes, each with a declared splitting structure over its unanswered attributes, on which a consistency-pruning rule tags inconsistent structure invalid without deleting it. The construction has two layers. The intrinsic layer takes a universe $U$ of such states, generated from a chosen set together with the blank state and the all-statements state and closed under the specified joins and joint assertions, and shows that $U$ is a commutative semiring, with informational union as addition and joint assertion as multiplication, under an explicit distributivity hypothesis that the unrestricted closure does not satisfy. The weighted layer builds the monoid semiring of measure-valued finite combinations over $U$ on the multiplicative monoid alone, so it requires no such hypothesis. Its principal result is that last-common-ancestor factorization is a derived identity of the bilinear multiplication rather than an imposed clause; through it every weighted sum admits the factored form. Admissible weight valuations are characterized as congruence-respecting normalized semiring valuations, distinct from empirical probability measures, and in models with finitely many attributes the doomed carriers are exactly the tagged ones and form a multiplicative ideal, so that the validity quotient is an exact semiring quotient. The pruning rule runs over proper childhood, the visible terms of a parent's declared decompositions, a weight-filtered subrelation of the content order, which, on the cone a fixed finite family of closure-independent carvings generates at real weights, is the join-semilattice of visible partial fillings of the parent's carved blanks. Declared splitting is treated as model structure under three stated compatibility assumptions, and the results that use it are conditional on them; the principal theorem sequence is stated for finite models with nonnegative real weights, with general negligibility relations a separated extension for which an explicit coefficient system is shown to satisfy every axiom together. A worked finite example carries the construction from carriers through declared splittings, proper children and doom to the reduced form and the validity quotient. The empirical probability assignment is not, in general, a semiring homomorphism for joint assertion. Which weighting it is, is the same question for every attribute, answered by the weighting scheme every blank carries; for quantum attributes that scheme is the Quantum Nexus kernel of the physical companion, whose weights are the Born weights.
\end{abstract}

\section{Introduction: the two faces}
\label{sec:intro}

The companion papers develop one object from two directions. The object is the observer-indexed \emph{Aethus}: a deductively closed partial assignment of answers to attributes, where an unanswered attribute --- a \emph{blank} --- is read not as ignorance about a hidden fact but as genuine superposition of the alternatives, and where a two-quantifier $\exists\forall$ consistency postulate---its quantifiers ranging over descendants at every depth of the reflexive, transitive refinement relation---prunes inconsistent structure --- \emph{pruning}, and the companions' \emph{deletion} idiom with it, being under this paper's tag architecture (§\ref{sec:reduced}) the realization-excluded reading of tags: carriers stand in the backend, and the frontend never unfolds into them. The modal face \cite{AethusSemantics} asks what it is to \emph{perturb} such an Aethus: reopen an entry, restate it, re-prune, and read a counterfactual off the residue. The present paper is the algebra face: it asks, and answers exactly, \emph{what algebraic structure the totality of such Aethae forms}, so that superposition, pruning, and weighting become operations in a single named object rather than devices invoked case by case. The two faces are complementary in the plain sense: every construction here has a modal reading there, and the perturbation semantics there presupposes precisely the closure and pruning structure that is here made algebra.

\emph{What is established.} Four things, at the grades the body states. The universe set $U$ of unweighted Aethae is a commutative semiring under the informational union and joint assertion, with distributivity priced as an explicit hypothesis, discharged for the classical case and shown to fail where agreeing alternatives are uncarved, (\S\ref{sec:intrinsic}). The weighted algebra $T$ over $U$ is identified as the monoid semiring $M[U]$, and last-common-ancestor factorization is derived as an identity of its bilinear multiplication, (\S\ref{sec:extrinsic}). On $T$ the proper-child and reduced-form apparatus is defined and the companions' third postulate is restated as ``elements without a reduced form are the invalid ones,'' with invalidity a label whose weights stand beneath it, (\S\ref{sec:reduced}). And the pruning rule's $\exists\forall$ form is shown equivalent to a fixed-point characterization, (\S\ref{sec:fixedpoint}). The no-additive-inverses constraint is then shown forced at the unweighted layer and priced at the weighted one, (\S\ref{sec:placement}).

The construction is organized in two layers, and the reason for the split is itself part of the content. The structural information carried by attributes --- (a) the ways in which jointly asserted attributes may be rewritten as alternative attributes encoding the same information, as in ``square'' versus ``rhombus $\cap$ rectangle,'' and (b) which specific intersections of attributes are \emph{labeled} inconsistent --- is foundational, and belongs to an unweighted, intrinsic layer. Weights are then a further, extrinsic degree of freedom layered over that foundation without changing the underlying carrier operations; they do affect which components count as proper children, and therefore which continuations the third postulate quantifies over, which is the intended feature. The first, ``unweighted/intrinsic'' layer accordingly constructs a universe set $U$ of unweighted Aethae and its commutative semiring structure, inheriting its foundations from the attribute apparatus of the fuller development \cite{Aethus}. The second, ``weighted/extrinsic'' layer builds the weighted algebra $T$ as a customizable measure-weighted superstructure over $U$, equipped with a factorization identity and a small quotient that jointly re-import the graph-theoretic content of the intrinsic layer back into its algebra.\footnote{The dating and origin note from the fuller development is retained here as provenance: the two-layer formalization was conceived on May 16, 2026 --- while the idea of splitting to two trees was being weighed, with an accompanying (non-load-bearing) mental image of a fleur-de-lis with upturned sides, the central shaft the unweighted layer $U$ and the full figure the weighted structure $T$ --- and written up for the long manuscript directly thereafter, as a proposed alternative method of defining proper child states from the supposition earlier in that paper.}

\section{The objects and maps, specified}\label{sec:spec}

\noindent This section fixes every domain and map the paper uses, with distinct notation wherever the domains differ; the sections that follow construct and justify them.

\begin{enumerate}
\item[(D1)] \emph{Carriers.} $U$, the Aethae of a working universe, each with its closed content $\mathrm{C}(A)$ (Definition~\ref{def:aethus}).
\item[(D2)] \emph{Intrinsic operations and order.} Joint assertion $A \cdot B$ and informational union $A + B$ on $U$, and the content order: $B \sqsubseteq A$ iff $\mathrm{C}(A) \subseteq \mathrm{C}(B)$, $B$ a \emph{child} of $A$ (Propositions~\ref{prop:order}, \ref{prop:lcd} and~\ref{prop:lca}).
\item[(D3)] \emph{Weighted combinations.} $T = M[U]$, the finitely supported $M$-valued functions on $U$, with the weighted-parent order $y \preceq_{\mathrm{add}} x$ (Lemma~\ref{lem:pointwise}).
\item[(D4)] \emph{Relative negligibility.} A relation $\mu \ll \nu$ on $M$ satisfying the visibility axioms of \S\ref{sec:extrinsic}; at the main domain $M = \mathbb{R}_{\ge 0}$, $\mu \ll \nu$ iff $\mu = 0 \ne \nu$.
\item[(D5)] \emph{Declared splitting data.} For each carrier $A$, its declared splittings $A \rightsquigarrow \sum_i w_i A_i$: finite families of cells, each a filling of a carving $A$ holds (Definition~\ref{def:descent}, rules (b) and (c)), with coefficients $w_i \in M$.
\item[(D6)] \emph{Proper childhood between carriers.} $B \le A$, with weight $B \le_\beta A$ (Definition~\ref{def:child-proper}): $B$ is a term of some admissible decomposition of $A$ whose coefficient is not negligible relative to the remainder.
\item[(D7)] \emph{Proper childhood in a combination.} $a \lhd_\omega x$, for a carrier $a \in U$ and a combination $x \in T$ (Definition~\ref{def:child}): $a$ carries in $x$ a coefficient not negligible relative to the rest of $x$. It relates a carrier to a combination, not two carriers, and is written differently for that reason.
\item[(D8)] \emph{The map between the two readings.} The embedding $\iota: U \to T$, $A \mapsto 1 \cdot A$. On embedded pure carriers the weighted validity criterion is exactly the unweighted one, although the embedding does not reproduce $A$'s declared proper descendants as weighted terms, the free-sum support of $\iota(A)$ being $A$ alone; for superpositions the correspondence runs in one direction, and in the other under cofinality (Proposition~\ref{prop:harmonize}).
\end{enumerate}

\textit{Declared splitting is model structure.} The free monoid semiring $T$ does not identify a carrier with its declared expansion: $A$ and $\sum_i w_i A_i$ are different elements of $T$. This paper does not construct an algebra in which they are identified. It treats declared splitting as additional model structure, subject to three compatibility assumptions. \emph{(S1) Composition}: if $A \rightsquigarrow \sum_i w_i A_i$ and $A_j \rightsquigarrow \sum_k v_k A_{jk}$ are declared, the composite $A \rightsquigarrow \sum_{i \ne j} w_i A_i + \sum_k w_j v_k A_{jk}$ is admissible. \emph{(S2) Coherence}: a cell's coefficient is the same in every admissible decomposition of $A$ in which it is a term. \emph{(S3) Additivity}: in a composed decomposition, the coefficients of the finest cells extending a cell sum to that cell's coefficient. The results that use splittings, the transitivity of proper childhood (Proposition~\ref{prop:child-transitive}), the join-semilattice of visible partial fillings (Theorem~\ref{thm:proper-lattice}) and every statement about reported weights, are conditional on (S1)--(S3), and the worked example of \S\ref{sec:model} satisfies them. An algebra that carries the splitting identities with well-defined operations would make them theorems; constructing it remains open, and nothing below is presented as supplying it.

\textit{The main domain.} The principal theorem sequence is stated for finite working universes and $M = \mathbb{R}_{\ge 0}$. Measure types with a nontrivial negligibility relation are a separated extension (\S\ref{sec:extrinsic}), for which Proposition~\ref{prop:polymodel} supplies a coefficient system satisfying every visibility axiom together.

\textit{What is proved about pruning.} Five results are distinct, and none supplies the others; Table~\ref{tab:pruning} gives each with its hypotheses. Throughout, an invalid carrier stays represented in $T$ under its tag; only its image in the validity quotient is zero.

\begin{table}[ht]
\centering
\small
\caption{Pruning and reduced-form results, with their hypotheses.}
\label{tab:pruning}
\begin{tabular}{@{}p{0.36\linewidth}p{0.34\linewidth}p{0.22\linewidth}@{}}
\toprule
Result & Hypotheses & Where \\
\midrule
Doom is hereditary along declared descent & Transitivity of $\le$, hence (S1) & Lemma~\ref{lem:fp} and the remark after it \\
Doom is a multiplicative ideal, $I = D$ & Finitely many attributes; declared descent by rules (a)--(c); closure's splitting rule read over proper children; tags carried by content, non-explosively & Proposition~\ref{prop:doom-tag} \\
A combination has a nonzero safe component iff some positively weighted carrier is undoomed & $M = \mathbb{R}_{\ge 0}$ & Proposition~\ref{prop:harmonize}~(ii) \\
A greatest reduced form exists, the valid-support truncation & $M = \mathbb{R}_{\ge 0}$; can fail in general, as with infinitesimal weights & Lemma~\ref{lem:basecase}; \S\ref{sec:extrinsic} \\
The validity quotient represents the reduced form & $M = \mathbb{R}_{\ge 0}$ and a valid blank carrier; can fail in general, as for $x = \varepsilon V + I$ & Lemma~\ref{lem:basecase}; \S\ref{sec:extrinsic} \\
\bottomrule
\end{tabular}
\end{table}

\section{Scaffolding: pins, states, Aethae, and the postulates}
\label{sec:scaffolding}

The minimum apparatus imported from the companions, stated once. An \emph{attribute} abstracts an event or property into an Aethus-relative possibility among alternatives --- a contingent truth customized for the informational setting --- ranging from Boolean predicates (``the right slit was traversed'') to dimensioned scalar quantities. Its semantic character is carried by the attribute itself; what the formalism captures is exactly the intrinsic compatibility structure across attributes, points (a) and (b) above, while deliberately shedding direct identification with the attribute's semantic instantiation in the world.\footnote{The stance is the one modern mathematics took with Euclid: retain the deductive apparatus of the \emph{Elements} while shedding the ontological characterizations of Book I's Definitions (``a point is that which has no part''); Hilbert's axiomatization makes the stance explicit by leaving ``point'' and ``line'' as undefined primitives characterized only by their relations under the axioms. The slogan: keep the structural content; let the ontological characterizations fall away.}

The objects are built from the companion's foundations \cite{Aethus}, restated here only as far as this paper consumes them. A \emph{pin} is a primitive atom carrying one primitive relation, disjointness; an \emph{attribute} is a set of at least two pairwise disjoint pins; a \emph{state} $[X : S]$ asserts that the attribute $X$ is filled by exactly the pins of $S$, determinate when $S$ is a single pin and blank otherwise, fully blank when $S$ is every pin of $X$ and semiblank when it is not; and a \emph{representation} is a set of states, read conjunctively. Deductive closure (§\ref{sec:order}) sends a representation to everything it commits to, and an Aethus is a class of representations with one closure (Definition~\ref{def:aethus}). An Aethus $a$ is a \emph{parent} of $b$ when some representation of $a$ is a subset of some representation of $b$ --- coarser information, higher in the tree. The \emph{refinement relation} $B \sqsubseteq A$ is closed-content inclusion, reflexive and transitive with descendants at every depth, and is stated independently of the weighted apparatus; the third postulate's relation, the physics-face companion's descent \cite{AethusQuant}, is the proper-child relation $\le$ of Definition~\ref{def:child-proper}, distinguished from $\sqsubseteq$ throughout.

\textit{The three decomposabilities, and the names of the grades.} An attribute $X$ is \emph{decomposed} by attributes $X_1, \ldots, X_n$ when it is their composite, its pins the joint pins of the components, so that $[X_1 : S_1] \cdot [X_2 : S_2] = [X_1 \wedge X_2 : S_1 \times S_2]$ under the joint assertion below. Relative to such a decomposition and to an Aethus, $X$ is \emph{state-decomposable} when every component is determinate, which holds exactly when $X$ is determinate; \emph{blank-decomposable} when every component is fully blank; and \emph{mixed-decomposable} when some components are determinate and the rest fully blank \cite{Aethus}. None of these is innate to the attribute: an attribute is decomposable in a given sense only relative to an Aethus, through its state form there, and the same attribute is state-decomposable in one Aethus, mixed-decomposable in a second and blank-decomposable in a third; what is innate to it is its pin set and its decompositions. The grades themselves are read off the filling of $X$ and are the same through every decomposition: \emph{blank} is anything non-determinate, and it divides into \emph{fully blank}, every pin of $X$ allowed, and \emph{semiblank}, two or more pins allowed but not all. A mixed-decomposable attribute is semiblank; a fully blank one is blank-decomposable through every decomposition, the converse failing only when the Aethus correlates the components, a correlation that lives on the composite alone. The companion's dichotomy theorem, in decomposition form, says that a mixed-decomposable attribute is the Aethic product of a determinate state, the composite of its determinate components, and a blank-decomposable one, the composite of its fully blank components \cite{Aethus}. This is what lets a state in an Aethus be referred to exactly by what its attribute decomposes into, ``$X$ with $X_1$ determinate at $a$ and $X_2$ fully blank,'' without writing the superposition's pins; the Aethic blank $\square$ of the statics below is the notation for a fully blank component. The intuition is the author's earliest for the framework, stated under the name \emph{attribute hybrid} in his journal entry of August 22, 2022. The procedure behind it is to picture a state as a product of other states and to let the written pins follow: the composition $[X_1 : S_1] \cdot [X_2 : S_2]$ above is that picture as a rule, and it is why a carving, which writes no set of pins in a state, is nonetheless a factor like any other, the carriers of the algebra being fillings, carvings and their joint assertions. The abstraction of products of states comes first, and a written set of pins is one representation of it.
% R889 (Fable 5.1, Oct 30 2026), at the author's direction: the product-first procedural intuition, with the carving as the factor that writes no pin set.
% R888 (Fable 5.1, Oct 30 2026), at the author's direction: decomposability stated as relative to an Aethus, never innate.
% R887 (Fable 5.1, Oct 30 2026), at the author's direction: the three decomposabilities, the names of the grades, and the dichotomy in decomposition form, imported from the long paper.
% R887 (Fable 5.1, Oct 30 2026), at the author's direction: blank made anything non-determinate, with the author's names fully blank and semiblank. Previous text: determinate when $S$ is a single pin and blank when it is every pin;

\textit{Child against proper child.} The two relations are defined once, with their transitivity and the comparative weight condition, in Definition~\ref{def:child-proper} and Proposition~\ref{prop:child-transitive}. Multiplicative childhood is the lattice descent $\sqsubseteq$ the semiring carries, reflexive and transitive ($A = A \cdot 1$, $A \cdot C \cdot D = A \cdot (C \cdot D)$), and the third postulate does not quantify over it: quantifying the postulate over the lattice relation would doom every Aethus, as the countermodel of Appendix~\ref{app:countermodels} shows. Neither relation has a single-level form, as \S\ref{sec:child-proper-child} states. Multiplicative childhood and declared proper childhood both relate carriers in $U$. The weighted-occurrence relation $a \lhd_\omega x$ instead relates a carrier to a combination in $T$. The following observation concerns that third relation, and it neither replaces the carrier-to-carrier relation nor establishes the representation correspondence of \S\ref{sec:spec}: \emph{if $a$ occurs as a proper-child carrier of $x \in T$, then $a$ is a child, in the multiplicative sense, of the informational core of $x$'s support}, $\mathrm{core}(x) := \sum_{A \in \mathrm{supp}(x)} A$, the last common ancestor of the carriers (in $U$ by closure under binary joins for finite support; for an infinite admissible support the join is an additional closure requirement the support class must state), with $A_i = \mathrm{core}(x) \cdot A_i$ since the core's content is contained in each $A_i$'s.

\textit{The two operations, and the inverted directions.} Two operations act at the representation level with directions inverted relative to their names, an inversion the fuller development derives and now proves as its axiomatic construction of the unweighted Aethus \cite{Aethus}: \emph{joint assertion}---the Aethic \emph{intersection}, $A \cdot B$---has content $\mathrm{Cl}(\mathrm{C}(A) \cup \mathrm{C}(B))$, everything either operand states, and is the \emph{least common descendant} of $A$ and $B$, their meet in the refinement order, lying in the child direction of both; and \emph{informational union}---the Aethic \emph{union}, $A + B$---has content $\mathrm{C}(A) \cap \mathrm{C}(B)$, what both operands state, and is the \emph{last common ancestor}, their join, lying in the parent direction. The intersection is never a strict ancestor and the union never a strict descendant, except in the equality cases; the set operation on content runs opposite to the name of the Aethic operation, which is why the names mislead a reader who takes them set-theoretically. And the two are constructed at different levels, which the equivalence-class definition of an Aethus forces. The intersection is built at the level of representations: $[S_A \cup S_B]$ is independent of the representations chosen, since $\mathrm{Cl}(S_A \cup S_B) = \mathrm{Cl}(\mathrm{Cl}\,S_A \cup \mathrm{Cl}\,S_B)$ for any closure operator, so joint assertion is the set union of representations, at any representations whatever. The union has no such construction: $S_A \cap S_B$ depends on which representations were chosen and is generically the wrong class, because representations need not be written with the attributes whose intersection is wanted, and no canonical choice repairs this. It is therefore defined over the abstract Aethae, in the order: the last common ancestor of $A$ and $B$ is the infimum under ancestry of the set of Aethae $x$ each of which constructs $A$ by intersecting with some $x_1$ and $B$ by intersecting with some $x_2$---$\inf_{\sqsupseteq}\{x : \exists x_1, x_2,\ A = x \cdot x_1,\ B = x \cdot x_2\}$---and the fuller development proves that this infimum exists, is unique, and has content $\mathrm{C}(A) \cap \mathrm{C}(B)$ \cite{Aethus}. The common content must then be read exactly, because the naive reading of it is wrong too. Content is \emph{closed} content: an item belongs to $A + B$ iff it is present in some representation of $A$ and in some representation of $B$, and representations include every item the Aethus's content derives, states on composites and exclusions among them. So $\mathrm{C}(A) \cap \mathrm{C}(B)$ is \emph{not} the intersection of the operands' atomic answers. Let the constraint layer contain $\{q, \neg t\}$ and $\{r, \neg t\}$, so that a further value $t$ is derivable from $q$ and from $r$; let $A$ state $p$ and $q$, and $B$ state $p$ and $r$, with $q$ and $r$ two values of one attribute $\alpha$. The atomic answers they share are $\{p\}$, but $t$ is derivable in both, hence present in both closed contents, hence in the union: $A + B$ states $p$ and $t$, and is a strict descendant of $[\{p\}]$. What the union does \emph{not} contain is an agreeing superposition of $\alpha$ over $\{q, r\}$. Closure never derives one from a single value, since each operand fills $\alpha$ with one value only; $A + B$, their last common ancestor, is instead the unweighted shadow of a \emph{disagreeing} superposition over its branches, and it carries the marks of one: $\alpha$'s single-pin carving, which both operands hold, and the exclusion of $\alpha$'s other values \cite{Aethus}; the disagreeing superposition proper, which records the branches, lives in the weighted layer. Agreeing superposition is asserted, never derived. Reading Aethae as sets of atomic answers ordered by inclusion, one is led to conclude that the last common ancestor is ``the intersection''; that gets the content wrong, since the union carries the restriction the atomic intersection loses, and the direction wrong, since the Aethic intersection is the descendant. The sentence ``informational union corresponds to set-theoretic intersection of stated-attribute sets'' invites exactly that reading, and is true only of the canonical closed representations. The restriction the union carries, the values of $\alpha$ that both operands exclude, is what distinguishes $A + B$ from $A + B'$ for a $B'$ stating $p$ and $s$: two Aethae the atomic intersection cannot tell apart.

\textit{Contradictions, and the label discipline.} A \emph{contradiction} is a jointness-disallowed combination within the intrinsic compatibility structure --- an attribute stated with two disjoint values, or a more elaborate circuit-level combination of attributes disallowed by their intrinsic jointness semantics --- and these are \emph{stipulated tags} rather than logical falsehoods: typically contingent in kind (``this apple is red all over'' against ``this apple is green all over,'' not $1 = 0$), perfectly consistent with the graph itself, passively marking the Aethae that carry them. These stipulated tags are the base cases from which the labeled invalidity of §\ref{sec:reduced} recursively propagates. Throughout, invalidity is operationally a \emph{label} assigned by the construction, not a directly realized empirical property; the connection between labeled invalidity and the empirical forbidding of an Aethus from being realized as an observer's actual informational Aethus is supplied separately, as a \emph{hypotheses non fingo} \cite{newton1687principia}--style injected realization principle, not derived within the formalism. This label-versus-realization discipline recurs at the section's close, where it fixes the logical form of the consistency postulate itself.

\section{Statics: the notation that denotes carriers}\label{sec:statics}

\begin{defin}[The statics-to-carriers map]\label{def:statics-map}
Let $\Phi$ send a draw static $S$ to the carrier $[\mathrm{Cl}(S)] \in U$ whose content is the closure of the fillings $S$ writes, and send each constraint entry $(g_1, \ldots, g_k) \mapsto \mathsf{F}$ of $S$ to the declaration in the relation structure that the joint assertion of the entries $g_1, \ldots, g_k$ is $0$, the tag the carrier carries wherever those entries are jointly realized in it. The map is the formal content of the sentence that a static denotes a carrier together with its tags and nothing else. \end{defin}

\begin{prop}[Compatibility of the statics-to-carriers map]\label{prop:statics-map}
(i) $\Phi$ is well defined on statics up to closure, two statics with the same closure denoting the same carrier and the same tags; (ii) $\Phi$ commutes with closure, $\Phi(\mathrm{Cl}\,S) = \Phi(S)$; (iii) $\Phi$ carries the joint assertion of two statics, the concatenation of their entry lists, to the product $\Phi(S) \cdot \Phi(S')$ of Proposition~\ref{prop:lcd}, the tags of the concatenation being those of either factor together with any declaration the concatenation newly realizes; (iv) a carving of $X$ into blocks, written as the family of constraint entries $\{[X : f]\} \mapsto \mathsf{F}$ for every $f$ meeting two or more blocks, is carried to the carving-induced family of tags, so that enforced splitting in the statics and in the algebra are one declaration read in two notations.
\end{prop}
\begin{proof}
(i) and (ii) are the long paper's closure proposition, closure being idempotent and the denotation factoring through it; (iii) is Proposition~\ref{prop:lcd} applied to the closures of the two entry lists, a declaration newly realized by the concatenation being one whose entries lie jointly in $\mathrm{Cl}(S \cup S')$ and in neither closure alone; (iv) is the enforced-splitting definition read in the statics' notation, each constraint entry of the family being one carving-induced tag.
\end{proof}

The carriers of the algebra below are fillings, carvings and their joint assertions, and the tags are the relation structure's joint invalidities; a carving is itself a carried family of such invalidities, each tagging the agreeing superposition over a filling that meets two of its blocks wherever that superposition is realized, held at its scope rather than everywhere, so that the rulebook has one kind of rule and a bare parent is one that carries no carving on the blank in question. The paper's written notation for them is the \emph{static}, the long paper's successor to the static Aethus: a written object that denotes a carrier of $U$ together with the tags written with it, and nothing else. It has two primitives, the \emph{Aethic blank} $\square$ and the \emph{Aethic false} $\mathsf{F}$, the second not logical falsehood but the mark of joint invalidity. A \emph{draw static} is a finite list of entries: \emph{attribute entries} $X \mapsto f$ with $f$ a nonempty set of $X$'s pins, $X \mapsto \square$ meaning the whole pin set and an absent entry read as fully blank, each writing the filling $[X : f]$; and \emph{constraint entries} $(g_1, \ldots, g_k) \mapsto \mathsf{F}$, each $g_j$ a pin or a filling, declaring that set jointly invalid, which is a tag. The carrier a static denotes is the joint assertion of its fillings, and its closure is the closure of that carrier with the tags it writes added to the relation structure. Every carrier holding at most one filling per attribute is written exactly by a draw static, and the denotation is not one-to-one, which is the notation's economy.
% R887 (Fable 5.1, Oct 30 2026), at the author's direction: matched to the vocabulary above. Previous text: $X \mapsto \square$ meaning the whole pin set and an absent entry read as blank,

The regime content of the declared-descent rules (Section~\ref{sec:model}, \textit{Regimes as content}) is written with the same two primitives. The carving $c_b$ of an attribute $X$ into blocks is the family of entries $[X : f] \mapsto \mathsf{F}$ for every filling $f$ meeting two or more blocks, the carving into single pins the case of singleton blocks; the irreducibility marker $b^{*}$ on a block is the family $[X : f] \mapsto \mathsf{F}$ for every nonempty proper sub-filling $f$ of it; and a blank with no such entry is the bare filling, which is regimes as content. The third postulate is what makes these entries operative: it assigns validity over declared descent, so that when $[X : f] \mapsto \mathsf{F}$ tags the proper child realizing the agreeing superposition over $f$, validity of the parent runs only through the children in which $X$ is carved, and the split is selected; without the postulate the same tags could stand in the relation structure and select nothing. A static $S$ is \emph{closed} when its carrier already holds every filling its closure derives, and $\mathrm{Cl}(S)$ is the draw static writing the fillings of the closure at their narrowest derived form together with every tag the closure derives; $\mathrm{Cl}(S)$ is closed, $S$ and $\mathrm{Cl}(S)$ have the same closure, and so a static denotes an Aethus, two statics denoting the same Aethus exactly when their closures coincide, by the long paper's closure proposition and fundamental theorem of Aethic equivalence; it thereby denotes an element of the validity quotient as well, which identifies further objects, distinct tagged carriers with different closures all mapping to zero, so that equality of closures is sufficient for equality in the quotient and not necessary. What remains to be supplied is a rewriting rule on entries that computes $\mathrm{Cl}(S)$ syntactically, the closure algorithm, and the correspondence between an implementation of it and the map $\Phi$ of Definition~\ref{def:statics-map}.

\section{The intrinsic layer: the unweighted semiring $U$}

\subsection{Child and proper child, as the framework defines them}\label{sec:child-proper-child}

Two relations on Aethae are in play throughout, and the framework fixes them additively and multiplicatively, in the author's own terms, before any tree or refinement language is used. We introduce child and proper child together because their distinction governs the subsequent construction. Childhood uses the intrinsic multiplication. Proper childhood additionally uses admissible weighted decompositions and relative negligibility, whose coefficient apparatus is specified in the weighted section below (\S\ref{sec:extrinsic}). Its introduction here anticipates that apparatus; it is not a construction from the unweighted semiring alone.

\begin{defin}[Child and proper child]\label{def:child-proper}
Let $A, B$ be Aethae with $A \ne 0$. $B$ is a \emph{child} of $A$ when $B = A \cdot C$ for some Aethus $C$: childhood is multiplicative, the joint assertion of $A$ with further content. $B$ is a \emph{component of $A$ with weight $\beta$} when $A$ admits an admissible declared decomposition \[ A \;\rightsquigarrow\; b\,B + c\,C, \qquad b \ne 0, \] written with $\rightsquigarrow$ because the expression $bB + cC$ belongs to $T$ but is not identified there with $1 \cdot A$ (\S\ref{sec:spec}), $\beta$ recording the coefficient $b$ against the remainder's weight $c$: $\beta = b/(b + c)$ where the measure type admits division (the base case $\mathbb{R}_{\ge 0}$, any semifield), the coefficient relative to the parent's total by the Relational Weight Principle, so that all components of a normalized decomposition carry weights summing to one. We call $\beta$ the \emph{relative component weight} of $B$ in $A$, and it is not by definition the probability with which $B$ is realized: invalid components carry raw weight but are never realized, so with $A \rightsquigarrow \tfrac12 B + \tfrac12 I$, $B$ valid and $I$ invalid, $B$'s relative component weight is $\tfrac12$ while its probability among the surviving outcomes is $1$. The probability of realized content is the weight in the validity quotient, the relative component weight renormalized over the valid components \cite{Aethus}; the physical probability assignment agrees with $\beta$ itself when every component is valid and the declared coefficients are the weighting scheme's own, and a model may declare structural coefficients that differ from the physical weights, as the physics-face companion does for structural visibility \cite{AethusQuant}, in which case the two are separate assignments; in a type without division, $\mathbb{N}$ say, the pair $(b, c)$ up to common rescaling by a non-zero-divisor scalar is the datum, no ratio being formed or needed, the comparison below running on coefficients. The reported coefficient is the maximum: within a fixed free-sum representation it is $B$'s normal-form coefficient, the largest amount of $B$ extractable from that sum, the long companion's supremum convention read there; that different admissible expansions of a carrier agree on it, the splittings not being identities of the free algebra, is a coherence obligation booked with the representation theorem below. The coefficients sum to the weight of $A$; $cC$ is shorthand for an arbitrary admissible remainder, a weighted combination the layer can add back, $c \ge 0$ its total weight with $c = 0$ allowed, so that the trivial declaration $A \rightsquigarrow 1 \cdot A + 0 \cdot C$ makes every Aethus its own proper child with weight $1$; $A \ne 0$ and $b \ne 0$ are required outright, since $0 \not\ll 0$ would otherwise admit zero-weight proper children. The \emph{boolean reduction} is the relation the third postulate and the Aethic tree consume, and its condition is comparative, the author's: $B$ is a \emph{proper child} of $A$, written $B \le A$ ($A \ge B$: $A$ is a proper parent of $B$), the order of the Aethic tree, iff \emph{some} admissible decomposition $A \rightsquigarrow bB + cC$ has its coefficient not measure-zero relative to the remainder, $b \not\ll c$ in the strict relative-negligibility relation $\ll$ of \S\ref{sec:extrinsic}, the relation defined existentially over the admissible decompositions, which compose, so that no closure is taken (where $b + c$ is invertible, scale covariance with $(b+c)^{-1}$ and back with $b + c$ gives $b \not\ll c$ iff $\beta \not\ll 1 - \beta$); and $B$ is a \emph{proper child of $A$ with weight $\beta$}, $B \le_\beta A$, iff $B \le A$ with component weight $\beta$. Proper childhood thus selects the components whose coefficients are not relatively negligible, their weights summing, where normalization is defined, to one less the negligible share.  At the base case $b \not\ll c$ is $\beta > 0$; in a graded type an area-grade component against a volume-grade remainder is no proper child, however valid in itself. A valid $B$ at relatively negligible coefficient remains an object, the postulate blind to that occurrence of it, its weight kept for the recovery lemma of \S\ref{sec:reduced}. The admissible decompositions are those the splitting structure declares (Definition~\ref{def:descent}, rule (b)).
\end{defin}

\begin{prop}[Transitivity, reflexivity, and the order]\label{prop:child-transitive} For admissible normalized decompositions whose coefficient system satisfies (M-comp), \eqref{eq:mcomp}, both relations are transitive; the boolean reduction of proper childhood, $\le$, is reflexive on nonzero Aethae and antisymmetric on carriers, hence a partial order there. \end{prop}
\begin{proof} For children, $B = AC$ and $D = BE$ give $D = A(CE)$ by associativity. For proper children, $A \rightsquigarrow bB + cC$ and $B \rightsquigarrow dD + eE$ with $b \not\ll c$ and $d \not\ll e$ compose, by (S1), to the declared decomposition \[ A \;\rightsquigarrow\; bd\,D + (be\,E + c\,C), \] its expression computed from the two by distributivity in $T$, with the remainder admissible by closure of the remainder construction; since proper childhood is existential over admissible decompositions, this composed decomposition witnesses $D \le A$ exactly when $bd \not\ll be + c$, against the \emph{combined} remainder, the reported coefficient entering only afterwards, $bd$ bounding it from below and $beE + cC$ its remainder from above where the coherence obligation of Definition~\ref{def:child-proper} is discharged (where normalization is defined the weights compose, $D$ having weight at least $\beta\delta$ in $A$, the chain rule); and that step neither positivity of products nor the listed properties of $\ll$ supply. It is secured by one stipulation on the measure type, the composition property the apparatus of \S\ref{sec:extrinsic} must satisfy over \emph{admissible normalized} coefficients, $b + c = 1$ and $d + e = 1$: \begin{equation}\label{eq:mcomp} \text{(M-comp)}\qquad b \not\ll c \ \text{ and } \ d \not\ll e \;\Longrightarrow\; bd \not\ll be + c . \end{equation} At the base case $\mathbb{R}_{\ge 0}$, $b \not\ll c$ is $b > 0$ and $d \not\ll e$ is $d > 0$, so $bd > 0$. In a graded type, with $\deg(\mu\nu) = \deg\mu + \deg\nu$, $\deg(\mu + \nu) = \max(\deg\mu, \deg\nu)$, $\deg 1 = 0$ and $\mu \ll \nu$ iff $\deg\mu < \deg\nu$, the normalization carries it: $\deg(b + c) = 0$ and $b \not\ll c$ force $\deg b = 0 \ge \deg c$, likewise $\deg d = 0 \ge \deg e$, so $\deg(bd) = 0 \ge \max(\deg e, \deg c) = \deg(be + c)$. Without the normalization, $\deg b = \deg c = 0$ with $\deg d = \deg e = -1$ satisfies the premises and violates the conclusion, so (M-comp) ranges over normalized coefficients only; the paper uses no type for which it fails. Reflexivity is the zero remainder. Antisymmetry is structural: by the splitting rule (a) of Definition~\ref{def:descent} a proper child distinct from its parent fills at least one blank and so has strictly larger closed content, and closed-content inclusion is antisymmetric, so two-way proper childhood forces identity. \end{proof}

\noindent The general lattice claim for the Aethic tree, an aspiration of the framework rather than a property established for $\le$, requires qualification. What can be proved, at the main domain and from the pins and the intrinsic operations of \S\ref{sec:intrinsic}, is this: the cone generated by a fixed finite family of closure-independent carvings is a join-semilattice, meets exist only for compatible fillings of positive weight, and a lattice is obtained only by adjoining an auxiliary bottom.

\begin{theorem}[Visible partial-filling cones: a join-semilattice and its bottom-adjoined completion]\label{thm:proper-lattice}
Work at the main domain $M = \mathbb{R}_{\ge 0}$. Let $A \ne 0$ be a carrier and $\mathcal{F}$ a fixed finite family of single-pin carvings that $A$ holds (rule (b) of Definition~\ref{def:descent}), of attributes $X_1, \dots, X_m$ with pins $P_1, \dots, P_m$, satisfying (S1)--(S3). A \emph{partial filling} is a partial function $f$ from $\{1, \dots, m\}$ to pins with $f(i) \in P_i$; its \emph{cell} is the carrier
\[
A_f \;:=\; A \cdot \prod_{i \in \mathrm{dom}\, f} [X_i : \{f(i)\}],
\]
with $A_\varnothing = A$; and its weight $w(f)$ is its coefficient in the composed decomposition of $A$ by the carvings of $\mathrm{dom}\, f$, the sum of the weights of the finest cells extending $f$ under (S3). Assume the carvings are \emph{closure-independent}: for partial fillings $f$ and $g$, $\mathrm{C}(A_f) \subseteq \mathrm{C}(A_g)$ iff $g$ extends $f$, so that $f \mapsto A_f$ is injective. Let $\mathcal{P}_A$ be the set of cells $A_f$ with $w(f) > 0$: the visible cells generated by composing carvings of $\mathcal{F}$, no other splitting being admitted in this construction. Then, under proper childhood:
\begin{enumerate}
\item[(i)] Every member of $\mathcal{P}_A$ is a proper child of $A$, and the proper children of $A$ obtained from $\mathcal{F}$ are exactly the members of $\mathcal{P}_A$.
\item[(ii)] For $A_f, A_g \in \mathcal{P}_A$: $A_g \le A_f$ iff $g$ extends $f$.
\item[(iii)] Any two members of $\mathcal{P}_A$ have a least upper bound in $\mathcal{P}_A$, the cell of their \emph{common sub-filling}, $A_f \vee A_g = A_{f \cap g}$, where $f \cap g$ is the restriction of either to the attributes on which both are defined and agree; two cells of one carving join at $A$. So $\mathcal{P}_A$ is a join-semilattice.
\item[(iv)] Two members have a greatest lower bound in $\mathcal{P}_A$ exactly when $f$ and $g$ are compatible, agreeing wherever both are defined, and $w(f \cup g) > 0$; it is then their joint assertion, $A_f \wedge A_g = A_{f \cup g} = A_f \cdot A_g$ (Proposition~\ref{prop:lcd}). Otherwise no member of $\mathcal{P}_A$ lies below both, so $\mathcal{P}_A$ is not in general a lattice.
\item[(v)] Adjoining a formal bottom $\bot$ below every member of $\mathcal{P}_A$ gives a lattice, the bottom-adjoined completion, in which two members with no common lower bound meet at $\bot$. The bottom is an order-theoretic auxiliary: it is not a carrier, not a proper child of $A$, and not a continuation over which the third postulate quantifies, and the all-statements carrier is not thereby made a proper child of anything (rule (a)).
\end{enumerate}
With all atom weights positive, $\mathcal{P}_A$ is the join-semilattice of all partial fillings ordered by reverse extension. $\mathcal{P}_A$ is the full proper-child cone of $A$ only when $\mathcal{F}$ exhausts the admissible splitting structure below $A$.
\end{theorem}
\begin{proof}
(i) A decomposition generated by $\mathcal{F}$ is a composition of its carvings, admissible by (S1), and its terms are the cells $A_f$, each with coefficient $w(f)$ by (S2)--(S3); at real weights a term is a proper child exactly when its coefficient is positive, its remainder being the rest of a total that is positive. Distinct partial fillings give distinct cells by closure-independence. Merging blocks gives no further cell: a carrier filling $X_i$ by two or more of its pins at once realizes the agreeing superposition the carving tags (rule (b) with Proposition~\ref{prop:statics-map}~(iv)), and arises from no declared splitting, by rule (a).

(ii) If $g$ extends $f$, then $A_g$ is a term of the composed decomposition of $A_f$ by the carvings of $\mathrm{dom}\, g \setminus \mathrm{dom}\, f$, which $A_f$ holds, a child inheriting its parent's declarations with its content. Its coefficient there is the local one, $w(g)/w(f)$, by the chain rule of (S2)--(S3): the parent-relative weight divided by $w(f) > 0$. That is positive since $w(g) > 0$, so $A_g \le A_f$. Conversely, $A_g \le A_f$ makes $A_g$ a refinement of $A_f$, so $\mathrm{C}(A_f) \subseteq \mathrm{C}(A_g)$, and $g$ extends $f$ by closure-independence.

(iii) Both $f$ and $g$ extend $f \cap g$, and $w(f \cap g) \ge w(f) > 0$, the finest cells extending $f$ being among those extending $f \cap g$, so $A_{f \cap g} \in \mathcal{P}_A$, and it lies above both by (ii). A member $A_h$ above both has $h \subseteq f$ and $h \subseteq g$ by (ii), so $h \subseteq f \cap g$ and $A_{f \cap g} \le A_h$. Two cells of one carving disagree on that attribute, so $f \cap g = \varnothing$ and their join is $A$.

(iv) For compatible $f$ and $g$ the joint assertion $A_f \cdot A_g$ is $A_{f \cup g}$, the meet of Proposition~\ref{prop:lcd}. If $w(f \cup g) > 0$ it lies in $\mathcal{P}_A$ and below both by (ii), and a member below both is $A_h$ with $h$ extending $f$ and $g$, hence $f \cup g$, so $A_h \le A_{f \cup g}$. If $f$ and $g$ assign different pins to one attribute, no partial filling extends both. If they are compatible but $w(f \cup g) = 0$, every $h$ extending $f \cup g$ has $w(h) \le w(f \cup g) = 0$, so no member of $\mathcal{P}_A$ lies below both.

(v) Adding $\bot$ below every element of a join-semilattice in which every pair with a common lower bound has a greatest one gives a lattice, the pairs with no common lower bound meeting at $\bot$. With all atom weights positive every cell is in $\mathcal{P}_A$, and $f \mapsto A_f$ is an order isomorphism from the partial fillings under reverse extension, by (i) and (ii).
\end{proof}

\noindent Remarks. The meet, where it exists, is the intrinsic meet; the join is the least cell above the intrinsic join of Proposition~\ref{prop:lca}, coinciding with it exactly when the two fillings share no consequence beyond their common sub-filling. Closure-independence is a genuine hypothesis. If the retained structure entails that two carved attributes take equal values, the fillings $X = 0$ and $Y = 0$ close to one carrier, so the cell map is not injective and the common sub-filling of the two fillings is empty although the carriers are equal; a constraint between carved attributes can break it the same way, since with binary $a$ and $b$ and $\{a = 1, b = 1\}$ jointly invalid, holding $a = 1$ entails $b = 0$, and the cells for $a1$ and for $a1b0$ coincide. Identifying the fillings that generate one closed carrier, and showing that the operations descend to the classes, is the general theorem this paper does not prove. Carvings acquired below a cell, by rule (c), lie outside the fixed family, and the theorem does not cover them. At a general measure type the statements need one further condition, that the coefficient of $A_g$ in $A_f$'s own decomposition and its parent-relative weight give the same visibility judgment against their remainders, which at real weights is division by $w(f) > 0$; this paper does not show that a nontrivial coefficient system supplies it, and the polynomial system of Proposition~\ref{prop:polymodel} cannot normalize nontrivially graded decompositions; its extension to rational functions, remarked after that proposition, can, though whether it meets the further compatibility condition is not shown. Blocks replace pins without change: for a held carving whose blocks are pin sets, the cells are the joint fillings by one block each, and what is excluded at every grain is the merged cell across blocks, which the carving tags. Two consequences govern the rest of the paper. First, proper childhood being transitive by composition, for admissible normalized decompositions under (M-comp), with no separate descent relation needed for arbitrary depth, the third postulate's ``two generations'' names its two nested quantifiers over the full relation $\le$: a validity floor must be supported by a proper child whose coefficient is not measure-zero relative to its parent's remainder, and a valid but relatively negligible component is neither a rescuing witness nor a relevant contrary continuation. Second, the tree and the postulate are formulated on proper childhood, multiplicative childhood supplying joint-assertion probes, the way forming $A \cdot C$ exposes an incompatibility. The free weighted algebra $T$ of \S\ref{sec:extrinsic} does not by itself encode the decompositions, $1 \cdot A$ having only the trivial one; the declared splittings on $U$ carry them, $T$ represents them within the scope of Proposition~\ref{prop:harmonize}, and the representation theorem that would close the gap is owed. The Relational Weight Principle fixes what such an identity can assert: no Aethus carries an absolute scale factor, the defined quantities are weight ratios and never bare coefficients, and rescaling every Aethus by one positive constant is an exact invariance, the weight structure being projective \cite{Aethus}. Within that, the splitting identities meet one obstruction:
\begin{prop}[A splitting obstruction]\label{prop:splitting-obstruction}
Let $B, C \in U$ be distinct valid carriers whose joint assertion is invalid, and let $b, c > 0$ at $M = \mathbb{R}_{\ge 0}$. In any quotient of the validity quotient in which carriers stay idempotent, the product stays bilinear, and coefficients on the two surviving carriers stay unique, $u[B] + v[C] = u'[B] + v'[C]$ only when $u = u'$ and $v = v'$, an identification $[A] = \lambda\,(b[B] + c[C])$ of a carrier $A$, for some $\lambda > 0$, forces $b = c = 1/\lambda$.
\end{prop}
\begin{proof}
In such a quotient $[B]^{2} = [B]$, $[C]^{2} = [C]$ and $[B][C] = 0$, so bilinearity gives $[A]^{2} = \lambda^{2} b^{2} [B] + \lambda^{2} c^{2} [C]$, while idempotence gives $[A]^{2} = [A] = \lambda b [B] + \lambda c [C]$. Uniqueness of coefficients on $[B]$ and $[C]$ then gives $\lambda^{2} b^{2} = \lambda b$ and $\lambda^{2} c^{2} = \lambda c$, whence $b = c = 1/\lambda$.
\end{proof}
\noindent\textit{What the obstruction proves.} Under the proposition's assumptions, the one identification of a carrier with a splitting of exclusive cells is the unit-coefficient sum $[A] = [B] + [C]$. The uniform mixture $\tfrac12[B] + \tfrac12[C]$ is excluded at its own scale, its square being $\tfrac14[B] + \tfrac14[C]$, but it has the same relative coefficients as $[B] + [C]$ on this fixed decomposition, the factor $\tfrac12$ being a bare coefficient the principle never treats as a defined quantity; identifying $[B] + [C]$ with $A$ remains part of the representation proposed below. Nonuniform splittings are obstructed even up to the principle's rescaling. The argument uses only carrier idempotence, the bilinear product, uniqueness of coefficients on the cells and their exclusivity, so a construction that closes the representation problem must say which of these changes: the embedding of carriers, the multiplication on weighted mixtures, the interpretation of splitting, or the semantics of the coefficients.

\textit{What a fixed exclusive family supplies.} One candidate, not developed here, changes the embedding of carriers and keeps the weights as coefficients: a carrier holding a carving is the unit-coefficient sum of its cells, idempotent in the validity quotient since $(B + C)^{2} = B + C + 2BC$ and $BC$ is invalid, while a nonuniformly weighted Aethus such as $\tfrac23 B + \tfrac13 C$, squaring to weights in ratio $4:1$ rather than $2:1$, is a distinct, non-idempotent element carrying the weights the carrier lacks. On a fixed family of exclusive valid cells with unique coefficients, $W = \sum_i w_i C_i$ with every $w_i > 0$ satisfies $W^{2} = \lambda W$ for some $\lambda > 0$ if and only if the $w_i$ are all equal, so uniform positive mixtures, and only they, admit an idempotent representative. This suggests reading the intrinsic algebra of \S\ref{sec:intrinsic}, from pins onward, as a decompositional algebra on such normalized support objects, each carrier decomposing, under a carving it holds, into its cells. Joint assertion is inherited, the product of two unit-coefficient sums being the unit-coefficient sum over the intersection of their supports. Addition is not: $B \sim 2B$ while $B + C \not\sim 2B + C$, the principle's invariance being one common rescaling of everything rather than an independent rescaling of the summands, so union must be the support union, a separate operation whose properties have to be proved.

\textit{What remains open across decompositions.} Constructing the intrinsic operations on these objects in general, and proving them compatible across decompositions, is open, and grain is the central difficulty: refining $B$ as $B_1 + B_2$ gives $bB = bB_1 + bB_2$, which doubles $B$'s mass if a coefficient is a mass, and is consistent as a density only relative to a reference measure with $m(B) = m(B_1) + m(B_2)$, a further construction the present algebra does not supply. Three readings of a coefficient are therefore kept apart: a mass, a density relative to a measure, and a relative weight within a declared decomposition. The established results use the last and rest on the declared-splitting formulation; the candidate carries none of them, and reconciling it with the preservation of a cell's share under refinement is the coherence obligation recorded below.

\subsection{The refinement order, proved once}
\label{sec:order}
This paper is the one the siblings cite for the statement that the graph of Aethae is well-defined, so the order-theoretic facts are proved here rather than described. Fix the pins, attributes, and states of §\ref{sec:scaffolding}, and let $\mathcal{L}$ be the set of \emph{items}: the states, together with the exclusions $X \not\ni p$ by which the companion's closure records that a pin no longer loads into the filling of its attribute---a set, so that the constructions stay in ordinary set theory. Let $\mathrm{Cl}$ be the companion's derivability closure \cite{Aethus}, generated by narrowing, propagation, elimination, invalidation, splitting over carvings, composition, and linkage. Invalidity stays external to it, in the companion as here: the empty filling that invalidation licenses stands as a tag on its carrier rather than exploding into every item, as the paragraph on carriers and Aethae below requires. $\mathrm{Cl}$ is extensive, monotone, idempotent, and \emph{non-weakening}: it never derives, from a single value of an attribute, a restriction of that attribute to a larger set of values, and every agreeing state it holds narrows one built by projection and combination from asserted states, a setting of a carving, or a state of an Aethus whose pin it holds, agreeing superposition being asserted rather than derived. A \emph{representation} is a subset $S \subseteq \mathcal{L}$, in practice a set of states, (a \emph{static Aethus} in the companion's original vocabulary, retained where its history is quoted); and the paper's object is defined once, its premise stated before its name.

\begin{defin}[Aethus]\label{def:aethus}
Call representations $S, S' \subseteq \mathcal{L}$ \emph{equivalent} when $\mathrm{Cl}(S) = \mathrm{Cl}(S')$. An equivalence class $A = [S]$ is an observer-indexed, deductively closed body of partial information: its \emph{content} is $\mathrm{C}(A) := \mathrm{Cl}(S)$ for any representation $S \in A$; its unstated attributes are not unknowns awaiting a hidden fact but constructive blanks, read as superposition of their alternatives; and it is indifferent to which of its representations presents it. An object with exactly these properties is what the framework calls an \emph{Aethus} (plural \emph{Aethae}). Validity is a predicate on Aethae (§\ref{sec:fp}) and not part of the definition of an Aethus.
\end{defin}

\noindent The premise is what earns the name. An Aethus is not a possible world, which settles every attribute; nor a set of worlds, which would enumerate instances the Aethus never collects; nor a theory whose unasserted sentences are merely undecided. Its blanks are superposition, and its identity is its closure, whatever representation presents it.

\noindent With the object named, the \emph{ambient Aethic lattice} is the powerset of the item universe quotiented into Aethae, $\mathcal{A} := \mathcal{P}(\mathcal{L})/\!\sim$, every combination of items into a representation taken up to derivability \cite{Aethus}. Every infimum and supremum in this subsection is taken in $\mathcal{A}$. The universe $U$ of Definition~\ref{def:universe} below is a different object: a \emph{working universe} $U_G \subseteq \mathcal{A}$ generated from a chosen set $G$ and closed under the operations the algebra needs, written $U$ when the generators are understood. The order and the two operations live on $\mathcal{A}$; a working universe inherits them, as the remark after Proposition~\ref{prop:lca} records.

\begin{prop}[The order is closed-content inclusion]
\label{prop:order}
For representations $S_A, S_B$: some representation of $[S_A]$ is a subset of some representation of $[S_B]$ iff $\mathrm{C}([S_A]) \subseteq \mathrm{C}([S_B])$. Writing $B \sqsubseteq A$ ($B$ a child of $A$, the more informative Aethus below) for $\mathrm{C}(A) \subseteq \mathrm{C}(B)$, the relation is well-defined on classes and is a partial order: reflexive, transitive, antisymmetric.
\end{prop}
\begin{proof}
If $S_A' \subseteq S_B'$ for representations of the two classes, monotonicity gives $\mathrm{Cl}(S_A') \subseteq \mathrm{Cl}(S_B')$, i.e., $\mathrm{C}(A) \subseteq \mathrm{C}(B)$. Conversely, the closed contents are themselves representations of their classes, and one is a subset of the other. The order properties are those of $\subseteq$ on closed sets, antisymmetry holding on classes because equal contents are one Aethus.
\end{proof}

\begin{prop}[Joint assertion is the least common descendant]
\label{prop:lcd}
For representations $S_A \in A$, $S_B \in B$, the class $[S_A \cup S_B]$ is independent of the representations chosen, since $\mathrm{Cl}(S \cup T) = \mathrm{Cl}(\mathrm{Cl}\,S \cup \mathrm{Cl}\,T)$ for any closure operator. Writing $A \cdot B := [S_A \cup S_B]$, $\mathrm{C}(A \cdot B) = \mathrm{Cl}(\mathrm{C}(A) \cup \mathrm{C}(B))$; $A \cdot B$ is a child of $A$ and of $B$, and every common child is a child of $A \cdot B$. It is the meet of $A$ and $B$, and it is never a strict ancestor of either, the equality cases being the only ancestry it has. \end{prop}

\noindent Under enforced splitting the product has two cases with one lattice position. Either no carried carving's declaration fires on $\mathrm{Cl}(\mathrm{C}(A) \cup \mathrm{C}(B))$ and $A \cdot B$ is the least common descendant with the tags $A$ and $B$ carried; or one fires, the union realizing an agreeing superposition over a filling that meets two or more blocks of a carried carving, and $A \cdot B$ is the same least common descendant now tagged, a distinct element of $U$ identified with the absorber only in the validity quotient. Validity is decided by the carried declarations, the enforced-splitting definition at work at the join.
\begin{proof}
$\mathrm{C}(A) \subseteq \mathrm{C}(A) \cup \mathrm{C}(B) \subseteq \mathrm{C}(A \cdot B)$; if $D \sqsubseteq A, B$ then $\mathrm{C}(A) \cup \mathrm{C}(B) \subseteq \mathrm{C}(D)$ and $\mathrm{C}(D)$ is closed, so $\mathrm{C}(A \cdot B) \subseteq \mathrm{C}(D)$. The last clause is antisymmetry.
\end{proof}

\begin{prop}[Informational union is the last common ancestor]
\label{prop:lca}
Define $A + B$ as the infimum under ancestry of the common ancestors of $A$ and $B$---the set of Aethae $x$ each of which constructs $A$ by joint assertion with some $x_1$ and $B$ with some $x_2$, $\{x : \exists x_1, x_2,\ A = x \cdot x_1,\ B = x \cdot x_2\}$. The infimum exists and is unique, and $\mathrm{C}(A + B) = \mathrm{C}(A) \cap \mathrm{C}(B)$.
\end{prop}
\begin{proof}
By Proposition~\ref{prop:lcd}, $A = x \cdot x_1$ for some $x_1$ iff $\mathrm{C}(x) \subseteq \mathrm{C}(A)$ (take $x_1 = A$ for the converse), so the set is $\{x : \mathrm{C}(x) \subseteq \mathrm{C}(A) \cap \mathrm{C}(B)\}$. The intersection of closed sets is closed, so $\mathrm{C}(A) \cap \mathrm{C}(B)$ is the content of an Aethus $L$; $L$ is in the set, and every member is an ancestor of $L$ by Proposition~\ref{prop:order}. So $L$ is the infimum, unique by antisymmetry.
\end{proof}

\noindent Two consequences fix the vocabulary. First, the union is constructed at the level of the abstract Aethae, not of representations: set-intersecting \emph{arbitrary} generating representations is presentation-dependent and generically wrong---$\{p,q\} \cap \{p,r\} = \{p\}$, while the representations $\{p,q,t\}$ and $\{p,r,t\}$ of the same two Aethae intersect to $\{p, t\}$---but intersecting the \emph{canonical} closed representations $\mathrm{C}(A)$ and $\mathrm{C}(B)$ is exactly right, and that is what Proposition~\ref{prop:lca} computes. Second, under ancestry the Aethae form a lattice, with $+$ the join and $\cdot$ the meet, the blank Aethus $\Null = [\varnothing]$ at the top and the all-statements carrier $[\mathcal{L}]$ at the bottom; and the lattice is not a tree. Two independent attributes give the diamond $\Null \sqsupseteq [\{p\}], [\{q\}] \sqsupseteq [\{p,q\}]$, in which $[\{p,q\}]$ has two incomparable parents. The formal object is therefore the \emph{Aethic refinement poset}---a lattice under the two operations---and ``the Aethic tree'' is the companion's name for the picture drawn from a chosen root, which this paper keeps as that and not as a claim about the Hasse diagram. And a working universe inherits the structure: if $U_G \subseteq \mathcal{A}$ is closed under parents and under joint assertion, then for $A, B \in U_G$ the ambient join $A + B$ lies in $U_G$, being a parent of $A$, and the ambient meet $A \cdot B$ lies in $U_G$ by product closure, so $U_G$ carries the restriction of the ambient order with the same meets and joins.

\paragraph{Two orientations, one order.} The refinement order can be read in either direction, and the names \emph{meet} and \emph{join} reverse with the reading, so the convention is fixed once:
\begin{center}\fbox{\parbox{0.92\linewidth}{\emph{Information order} $A \le_I B \iff \mathrm{C}(A) \subseteq \mathrm{C}(B)$ ($B$ at least as informative as $A$; the blank Aethus at the bottom). \emph{Ancestry order} $B \sqsubseteq A \iff A \le_I B$ ($B$ a descendant of $A$; the blank Aethus at the top). Under $\le_I$: $A + B = A \wedge_I B$ (meet) and $A \cdot B = A \vee_I B$ (join). Under ancestry the names reverse: $A + B$ is the join (last common ancestor) and $A \cdot B$ the meet (least common descendant). Every ``meet'' and ``join'' in this paper names its order.}}\end{center}
\noindent Propositions~\ref{prop:lcd} and \ref{prop:lca} are stated in the ancestry orientation, which is the framework's; the information orientation is the conventional one in order theory, and a reader who prefers it loses nothing but must swap the two names.

\paragraph{Agreement belongs to the intrinsic layer.} One typing remark protects a distinction the whole corpus relies on. Whether two alternatives of a blank stand in \emph{agreeing} or \emph{disagreeing} superposition---whether each states the other's negation in its interior, the distinction the composition companion owns \cite{AethusUnion}---is a predicate on the intrinsic layer $U$: it is fixed by which states are present, the single-pin carving marking the disagreeing case \cite{Aethus}, and it is definable before the weighted algebra $M[U]$, before weights, path kernels, or phases exist. Coherence and decoherence are not: they are quantitative properties of a physical realization of that structure, available only once a kernel is installed on $M[U]$, and no definitional biconditional between the two vocabularies is stated anywhere in this corpus. Agreeing/disagreeing is the unweighted logical category; coherent/decoherent is its kernel-level physics; replacing the one word by the other is a change of level, not a synonym.

\paragraph{Carriers and Aethae.} $U$ contains, by closure under joint assertion, carriers that no observer's information Aethus could be: tagged jointly disallowed combinations, and the all-statements zero. Write $\widehat{U}$ for the universe of formal Aethus carriers just constructed and $U_{\mathrm{val}} \subseteq \widehat{U}$ for those not tagged invalid; ``information Aethus'' in the realizable sense means an element of $U_{\mathrm{val}}$, and the algebra runs on $\widehat{U}$ because closure under joint assertion requires the tagged carriers to stand. Where the text below says $U$ without qualification it means $\widehat{U}$. The companion's closure has no explosion rule either: there too the empty filling is a tag, so every carrier is an Aethus in the companion's sense, the tagged carriers included and kept distinct, and the coarse grain is the companion's validity quotient, which identifies every tagged carrier with the additive identity and on which probability is valued \cite{Aethus}. The carriers are distinct because the weighted layer weights them separately, as the reduced form of §\ref{sec:reduced} and the weighted reopened blank require.

\label{sec:intrinsic}

We construct an arbitrary ``universe set'' of unweighted Aethae --- implicitly, by conditions, rather than by containing all possible Aethae, much as a set-theoretic methodology has a universal set without enforcing that all possible sets be contained in it.

\noindent\textit{Two relations and one representation, named.} Throughout, \emph{refinement}, also called \emph{lattice descent}, is the closed-content order $\sqsubseteq$ of Proposition~\ref{prop:order}, the order of meets and joins, and is childhood itself, $B \sqsubseteq A$ iff $B = A \cdot C$ for some $C$; \emph{proper childhood}, Definition~\ref{def:child-proper}, over which the third postulate quantifies, is written $B \le A$, the order of the tree, $B \le_\beta A$ carrying the weight: a proper child is a refinement of its parent, each term of a union carrying the union's content, and proper childhood is the subrelation of refinement made of the visible terms of the parent's decompositions, \emph{declared descent} where a model's declared splittings are in view; and \emph{weighted expansion} is the representation of proper childhood that the free algebra $T$ affords, the expansion of a carrier into its weighted terms, within the scope Proposition~\ref{prop:harmonize} states, proposed beyond it pending the representation theorem of \S\ref{sec:child-proper-child}. Where ``descent'' stands unqualified below it means refinement.

\begin{defin}[Declared descent: the third postulate's relation]\label{def:descent}
The third postulate looks down proper childhood, the refinements of an Aethus that are visible terms of its declared splittings; its relation is \emph{declared descent}. A model of the algebra specifies, for each Aethus $A$ of its working universe, the set of $A$'s \emph{declared proper children}, the Aethae $A$ admits as its own splittings, subject to three rules.
\begin{enumerate}
\item[(a)] Every carrier is its own proper child at full weight, the reflexive case. A nontrivial declared splitting of $A$ yields descendants of $A$ under $\sqsubseteq$ that resolve blanks of $A$ along its standing structure, one value each, introducing no rival assignment arbitrarily; a resulting carrier may still be tagged invalid, its resolved content conflicting with retained constraints, and the tag does not remove it from the relation, such tagged continuations being what the postulate inspects. Carriers holding rival determinate values of one attribute, the all-statements Aethus among them, arise from no declared splitting and are proper children of nothing but themselves.
\item[(b)] An attribute whose single-pin carving $A$ holds, so that its blank is the union of the attribute's determinate states---an \emph{exclusive} blank, its alternatives excluding one another---splits: each filling of it by one alternative, closed, is a \emph{component} of $A$, a term of the decomposition the carving supplies, $A \rightsquigarrow \sum_i w_i\,(A \cdot [X : \{a_i\}])$ with the cells' parent-conditional weights $w_i$ wherever the model is weighted, and a \emph{declared proper child} of $A$ exactly when its coefficient is not measure-zero relative to the rest of that decomposition (Definition~\ref{def:child-proper}); in the unweighted model every cell is visible and the components are the proper children. A filling asserted without the carving puts the attribute's open pins on the table and declares no split, saying nothing about validity below it; an attribute $A$ does not state at all is such a filling, its pins all open.
\item[(c)] An attribute $A$ holds in \emph{irreducible} agreeing superposition---holding the companion's pin $e_S$, that the attribute is filled by exactly $S$ and admits no narrower filling \cite{Aethus}---splits only when $A$ itself carries the access to the split, a joint assertion with a carrier that records the attribute's value; a record of the value holds the attribute's carving, so the split is the carving's, as in (b), and every such split lies in $\mathrm{Base}$, a narrower filling being jointly invalid with $e_S$ by the companion's invalidation rule. An attribute held in agreeing superposition only in the broad sense, without $e_S$, likewise splits once access is granted, by the carving the record carries, and its splittings may then be valid; until then it is a filling, and declares no split.
\end{enumerate}
The declared splittings and their compositions are the admissible decompositions over which proper childhood $\le$ is defined existentially (Definition~\ref{def:child-proper}); no closure is taken, the relation being transitive by composition and no step size privileged. Which splittings an Aethus admits is a fact about that Aethus, and part of the data a model supplies, not a consequence of the algebra: an Aethus can carry the constraint of an irreducible agreeing superposition and still declare proper children that split it, which rule (c) then tags, or declare none. Which of the two an Aethus of nature is, is read off the situation modelled, in the absence of a detector from whether anything gives access to the split, and that reading belongs to the physics-face companion \cite{AethusQuant}, not to the algebra.
\end{defin}

\begin{prop}[Declared descent is weight-filtered refinement; how declared children join]\label{prop:descent} Let $A$ be an Aethus stating at most one value per attribute, and let $X$ be a blank of $A$ with pins $\{p\} \cup R$. \begin{enumerate} \item[(i)] \emph{Declared descent.} $B \le A$ iff $B$ is a descendant of $A$ under $\sqsubseteq$ reached by filling some of $A$'s exclusive blanks with one value each along carvings $A$ holds, each cell a component at a coefficient not measure-zero relative to the rest of its carving's decomposition; a bare blank, one whose carving $A$ does not hold, has no declared children under rule (b), and a relatively negligible cell is no declared child though it is a descendant under $\sqsubseteq$.
\item[(ii)] \emph{Carved parent.} If $A$ holds $c_X$, any two determinate children $A \cdot [X : \{p\}]$ and $A \cdot [X : \{q\}]$ join to $A$ together with the exclusions of the pins neither holds, and the determinate children over all the pins of $X$ join to $A$ itself.
\item[(iii)] \emph{Bare parent.} If $A$ does not hold $c_X$, the products $A \cdot [X : \{a\}]$ are carriers but not declared children; any two of them join to $A \cdot c_X$ together with the exclusions of the pins neither holds, and the family over all the pins joins to $A \cdot c_X$, the carving supplied by the determinate states through closure and not held by $A$, so the join holds a carving the parent lacks; and $A \cdot [X : \{p\}]$ joins back to $A$ itself with the filling $A \cdot [X : R]$ when $R$ holds at least two pins, that filling being a valid carrier of the bare parent and tagged under the carved one.
\end{enumerate}
\end{prop}
\begin{proof}
The first claim is rules (a) and (b) read together with transitivity: a chain of single-value fillings is a single-value filling, and conversely a filling of several exclusive blanks is reached by filling them one at a time. For the second, by Proposition~\ref{prop:lca} a join holds exactly what its operands' closed contents share. A determinate filling holds the single-pin carving by semantic mining and excludes every pin it omits \cite{Aethus}, so two determinate fillings share $c_X$ and the exclusions of the pins neither holds, and the family over all the pins shares $c_X$ alone, each exclusion being omitted by the filling that holds the excluded pin. In case (ii) $A$ already holds $c_X$, so what the children share beyond $A$'s own content is the stated exclusions and the full family shares exactly $A$; in case (iii) $A$ lacks $c_X$, so the shared carving is an addition and the joins are $A \cdot c_X$ with the same exclusions. The filling $[X : R]$ with $|R| \ge 2$ holds no carving of $X$ and excludes only $p$, while $[X : \{p\}]$ excludes only the pins of $R$, and no state of $X$ is common, the closure deriving no wider state from either; so $\mathrm{C}(B) \cap \mathrm{C}(C) = \mathrm{C}(A)$. The completed model of §\ref{sec:model} computes all three cases.
\end{proof}

\noindent So declared descent orders the single-valued Aethae as refinement, but joining its children does not undo it: the join of disagreeing determinate children is the core of the disagreeing superposition over them. The proposition keeps the two parents apart in its clauses (ii) and (iii): only in the bare case does the join hold a carving the blank lacks. That is the unweighted face of the difference between a blank and the disagreeing superposition over it, which the weighted algebra keeps as terms. What the declaration adds is the treatment of agreeing superposition, where the two readings of the companion's chess idiom part: an irreducible agreeing superposition has no declared children beneath it, and is judged on its own tag; an agreeing superposition that has acquired access to its split has children beneath it, all tagged, and is doomed by them, and, holding the carving its record carries with every setting tagged, it is tagged by closure's splitting rule as well.

\begin{defin}[Universe set of unweighted Aethae]\label{def:universe}
A universe set $U$ is constructed in two steps.
\begin{enumerate}
\item Construct an arbitrary \textbf{generator set} $G$ of unweighted Aethae. Every pair $a, b \in G$ has a \emph{last common ancestor} in the refinement order---namely their informational union itself---and this is not a premise but a theorem of the fuller development's axiomatic construction, the union being the infimum under ancestry of the common ancestors, which exists and is unique for every pair \cite{Aethus}. The trivial cases: $a = b$, whose last common ancestor is itself; and the pair sharing no items in any representations---a condition rarely met, since two Aethae determinate on one attribute at different pins already share the exclusion of its other pins and its single-pin carving \cite{Aethus}, so that their last common ancestor is the union carrying these and not the blank Aethus---whose last common ancestor is the canonical \textbf{blank Aethus} $\Null$, formally the Aethus whose representation is empty, itself a primitive unweighted Aethus. Because parentage is defined through representations (a subset relation on representations), the set of parents of each $a \in G$ is fully deductively intrinsic to $a$; the union premises depend on $G$ only through which branch points of those intrinsic trees happen to be referenced by pairs in $G$.
\item Construct the \textbf{universe set} $U$ as the closure of $G \cup \{\Null, \Hyp\}$ under four conditions, the ambient blank Aethus $\Null$ and the all-statements Aethus $\Hyp$ being placed in the generating set before closure, since an arbitrary $G$ need not generate either: $G \cup \{\Null, \Hyp\} \subset U$; for any $a, b \in U$, their joint assertion and their informational union are in $U$ --- tagged combinations included, standing as elements; and the ambient all-statements Aethus $\Hyp$ is adjoined as an element, the absorbing element of $U$, written $\bar 0$ where its algebraic role rather than its content is in view: not a new formal symbol but the all-statements carrier itself. Three other zeros live at other layers and are not this one: the weighted algebra's additive zero $0_T$ is the empty combination, no visible components, and is distinct from $\Hyp$ taken as a carrier of visible weight; the validity quotient's zero is the class of the tagged carriers under the validity congruence, into which $\Hyp$ and every tagged carrier map; and the empty filling that tags a carrier is content, not an element of $U$ at all. Parent-closure is not required, and it cannot be: every ambient Aethus is a parent of the all-statements Aethus, so parent-closure together with its adjunction would collapse $U_G$ into the whole ambient lattice. The finite models compute $U_G$ exactly as stated, as the closure of the generators under products and joins. Once adjoined, both semiring identities hold: $x + \Hyp = x$ (the all-statements representation intersects with any representation of $x$ to a representation of $x$) and $x \cdot \Hyp = \Hyp$ (by construction); and, being the all-statements Aethus, $\Hyp$ is itself tagged --- harmlessly, its role as identity being algebraic rather than validity-theoretic.
\end{enumerate}
$U$ is then given the structure of a \textbf{commutative semiring under Hypothesis~(D) below} --- unconditionally, a pair of commutative monoids obeying the two absorption laws, the bounded-lattice pattern, with the \emph{semiring} name earned exactly where (D) holds; $T$'s semiring structure is by contrast unconditional, consuming only $U$'s multiplicative monoid (\S\ref{sec:placement2}), so that (D)'s entire footprint stays the single consumer named there. The operations: addition is the informational union; multiplication is the joint assertion \emph{simpliciter}, jointness-disallowed combinations included --- these stand in $U$ as tagged elements rather than collapsing, invalidity riding as a label on the standing carrier; the additive identity is the \textbf{absorbing element} $\Hyp$ (multiplicatively absorbing by construction: joint assertion with the all-statements Aethus returns it), the multiplicative identity the blank Aethus $\Null$.
\end{defin}

Two structural remarks, in place of a declaration. First, the closure condition on joint assertions is \emph{necessary}, and an earlier shortcut fails without it: if out-of-universe products are instead piecewise replaced by $\Hyp$, distributivity breaks --- with $G = \{a, b, c\}$ pairwise sharing no items, $a \cdot (b + c) = a \cdot \Null = a$ while $a \cdot b = a \cdot c = \Hyp$ gives the right side $\Hyp$, and the failure is purely the not-in-$U$ replacement, no semantic invalidity involved. This paper accordingly retains one definition of multiplication (the joint assertion) and drops the source's second, factorization-based formulation, whose equivalence to the first was asserted rather than proven. Second, with the closure in place, distributivity of joint assertion over the coarsening join is stated as an explicit hypothesis:

\begin{quote}
\textbf{Hypothesis (D).} On the closed universe $U$, joint assertion distributes over informational union: $a \cdot (b + c) = a\cdot b + a \cdot c$. The hypothesis has a sufficient condition that is a theorem schema rather than a posit: \emph{if the consequence relation's Lindenbaum lattice is distributive and the invalidity tags are external to closure, then the lattice of closed contents ordered by inclusion is distributive, hence so is $U$ under $\le_I$, hence (D) holds}---the standard fact that the closed sets of a closure operator arising from a distributive consequence relation form a distributive lattice. The positive classical fragment has a distributive Lindenbaum lattice, so (D) holds there; for the framework's own closure the hypothesis is retained as a hypothesis because it holds only over carved unions, as the close of this discussion shows.
\end{quote}

\noindent The informational union is what the companion now calls the \emph{core} of the Aethic union \cite{Aethus}. There the union keeps its operands as terms, being the support of a weighted sum of this paper's kind with its weights forgotten, and joint assertion distributes over it exactly; the informational union is the content its terms share. Hypothesis (D) is accordingly a claim about cores, and where it fails, at agreeing alternatives that no carving separates, the failure is the core's and not the composition's.

\noindent For attribute systems governed by classical consequence, (D) is the distributivity of the filter lattice of a Boolean algebra and holds; and the clash-grade ruling simplifies (D)'s status outright --- with closure not exploding on tagged content, though it consults tags in the splitting rule, invalidity no longer bears on either operation, so the discharge for classical consequence reduces to the pure theory-lattice question, the non-Boolean case being the only remaining price --- with one scope clause the tag-blindness commitment forces: ``classical consequence'' here means the disjointness-excised (positive) fragment, since were value-disjointness axiomatized in the consequence relation itself, $\{q = v_1,\ q = v_2\}$ would explode classically and the standing-carriers commitment would fail; the discharge path survives, the Lindenbaum algebra of the positive fragment being a distributive lattice.

\textit{The status of (D).} (D) is accordingly priced here as a hypothesis at the papers' rigor grade, with its discharge under stated conditions a named obligation.

\textit{The classical discharge.} The chain, with the scope limitation below: with Aethae as deductively closed theories, the coarsening join is theory intersection and joint assertion is closure of the union, so the two operations on $U$ are the meet and join of the filter lattice of the Lindenbaum algebra; the filter lattice of a distributive lattice is distributive, and $U$ inherits (sums stay in $U$ by the binary-join closure the construction specifies, $a + b$ being the join of $a$ and $b$; products by the closure clause) --- so (D) would hold whenever the consequence relation's Lindenbaum algebra is a distributive lattice, the positive classical and intuitionistic fragments included --- the route presupposing a lattice fragment, $\wedge$ \emph{and} $\vee$ present: vocabularies lacking disjunction, including the flat vocabularies the trivial-LCA case elsewhere presumes, are not covered by it, falling to the residual or needing a direct theory-lattice argument, bare closure-system lattices being non-distributive in general --- leaving as residual the \emph{non-distributive} theory lattices. (The distributivity counterexample leans flat while this discharge leans rich; the generator sets differ, and there is no conflict.) The framework's own closure is not covered by this chain, and (D) does not hold for it without restriction. Agreeing superpositions are Aethus \emph{content}, jointly held alternatives \cite{AethusUnion}, and content of that kind blocks case analysis: with $[X : \{b, c\}]$ held and each of $b$ and $c$ entailing $z$, the products with $[X : \{a, b\}]$ and with $[X : \{c\}]$ both derive $z$, while the product with their union, which holds nothing on $X$, does not. What holds is the companion's theorem that joint assertion distributes over every carved union---over the settings of any carving, and in particular over every union of determinate Aethae, which keeps the single-pin carving they share \cite{Aethus}. That is the classical regime, where records carve; (D) fails only for unions of agreeing alternatives that no carving separates, the framework's located form of the non-distributive join of quantum logic.

\textit{Associativity without (D); invalid inheritance kept.} One compensation falls out without (D): \emph{associativity} of the multiplication, now immediate --- with tagged combinations standing rather than collapsing, every partial product lies in $U$ by closure and all groupings agree by the intrinsic associativity of joint assertion (set-theoretic union at the representation level). The tag apparatus keeps \emph{invalid inheritance}---every declared descendant of an invalid Aethus invalid, and every product with an invalid factor invalid. The first is not an axiom but a corollary of the third postulate on declared descent, in three lines: with \emph{formal invalidity} defined outright as membership in the least base-generated doom set, $D := G(\mathrm{Base})$ of §\ref{sec:fp} over the declared descent of Definition~\ref{def:descent}---which the transitive-descent lemma there shows is already the least fixed point---and formal validity as its complement, $\mathcal{V} := U \setminus D$: if $A \in D$ then every declared descendant $B \le A$ has every declared descendant of $B$ among the declared descendants of $A$, so every declared descendant of $B$ has a base case beneath it and $B \in D$. The second is the ideal property recorded in §\ref{sec:fp}: the quotient kills the monoid ideal the doomed set generates, which coincides with the doomed set once each attribute's regime is stated as content (Proposition~\ref{prop:doom-tag}). The definition settles a question the postulate alone leaves open, whether validity is the \emph{greatest} predicate satisfying its condition or merely some post-fixed point of it: the postulate as stated gives only $\mathcal{V} \subseteq F(\mathcal{V})$, from which $A \in F(\mathcal{V})$ does not yield $A \in \mathcal{V}$, and inheritance needs the identification $\mathcal{V} = \nu F$, which the definition supplies by fiat. What the definition does \emph{not} supply is empirical realizability: it fixes the presence or absence of a formal tag, and the passage from the tag to the physical forbidding of an Aethus remains the separately injected realization principle. Monotonicity along $\sqsubseteq$ and its chains follows, and the property now serves the propagation machinery rather than the multiplication's well-definedness. (Two prior justifications are retired: a parent-closure argument native to the abandoned piecewise multiplication, and a monotonicity route needed only while labeled products collapsed to $\Hyp$.)

The names $\Null$ and $\Hyp$ are this paper's neutral renderings of the fuller development's Null and hyperinvalid Aethae \cite{Aethus}; nothing beyond the identities they serve is imported with the names.

\section{The extrinsic layer: the weighted algebra $T$}
\label{sec:extrinsic}

\subsection{The weighted superstructure and its measure type}
\noindent\textit{Reading guide for this section.} The main results of the paper are stated, wherever a residual would otherwise enter, on the base case $M = \mathbb{R}_{\ge 0}$---or at most a cancellative, zerosumfree commutative coefficient semiring---and that base case is the stable core a reader should audit first: normal forms, coefficientwise order, visible weights, and the valid-support truncation as the reduced form. The generalization to arbitrary measure types with a relative-negligibility relation, developed from the paragraph \emph{Visibility, relative} onward and consumed by the modulo reading of §\ref{sec:reduced}, is the framework's extension and is quarantined from the core: none of the core theorems depends on it, its own open residuals are listed in §\ref{sec:placement}, two of which are already failures of the unrestricted generalization rather than unverified claims: with $x = \varepsilon V + I$, $V$ valid, $I$ invalid and $\varepsilon \ll 1$, the anchored reduced-form set is empty while $[x] = \varepsilon [V] \ne 0$ in the contracted semiring, so postulate-invalidity under relative visibility does not coincide with the quotient's zero class; and with positive hyperreal weights and infinitesimal negligibility, $y_\varepsilon = V + \varepsilon I$ is proper for every positive infinitesimal $\varepsilon$ with no largest one, so a maximum need not exist, the real-weight truncation theorem surviving both and being the one kept central, and a reader who takes $M = \mathbb{R}_{\ge 0}$ throughout loses no theorem of the paper.


To a universe set we associate a weighted algebra $T$,\footnote{Forward-looking note for any future categorical translation, retained from the source: morphisms between weighted algebras would most naturally be structure-preserving maps on the underlying unweighted layer $U$, with the externally imposed weight assignments from $M$ transported covariantly rather than treated as fixed structure to preserve. The measure structure is peripheral in the precise sense that it does not exist on $G$ or $U$ themselves but is externally imposed only in constructing $T$; correspondingly it is not validity-preserving, validity being a label assigned pre-realization rather than a structural property preserved under structural maps.} taken as bijectively correspondent to the tuple of a universe set $U$ and a measure-theoretic measure type $M$ (for instance, the Lebesgue measure restricted to its finite values --- the cancellativity stipulation below confines $M$ to the finite cone, $\infty + 1 = \infty + 2$ being the standard counterexample), $M$ being the formal generalization of real-valued weights, stipulated a \textbf{commutative semiring with unit} --- the products $\mu\nu$ serving $T$'s multiplication, the unit $1$ embedding $U \hookrightarrow T$ as weight-one terms (Proposition~\ref{prop:factor}'s left factor is $1 \cdot C$), and $M$-distributivity with $0$-annihilation serving the congruence discharge below --- and equipped, on its additive side, with two further stipulations and one definition --- the positive cone, additionally, being derived rather than imposed: under the algebraic order below, $0 \leq \mu$ for every $\mu$, so the cone is all of $M$ and ``non-negative'' is automatic wherever the constructions below invoke it.

\textit{The algebraic order, cancellativity, and zerosumfreeness.} The \emph{algebraic} (difference) order on $M$ --- a definition rather than an axiom: $\nu \leq \mu$ iff $\mu = \nu + \lambda$ for some $\lambda \in M$. Addition stipulated \emph{cancellative}, so that the residues the definitions below consume ($w - \mu$ in proper childhood; coefficient-wise differences in weighted parenthood) not only exist within the order but are unique, licensing the difference notation --- cancellativity being what residue \emph{uniqueness} consumes --- coefficient well-definedness itself needs no cancellativity, falling out of the monoid-semiring identification below --- and compatible with the non-Archimedean ambitions. And \emph{zerosumfreeness} --- no two nonzero elements summing to zero --- the guard §\ref{sec:placement} identifies: with cancellativity in place, antisymmetry of the algebraic order is \emph{equivalent} to zerosumfreeness, so that, under the standing additive-cancellativity assumption, this stipulation is what makes the algebraic preorder a partial order and earns the order vocabulary, while forbidding nontrivial inverses outright.

\textit{Visibility, relative.} Alongside $M$ itself, visibility is axiomatized, and it is \emph{relative}: the third postulate ignores anything measure-zero relative to anything else there. The formalization: a \textbf{negligibility relation} $\mu \ll \nu$ (``$\mu$ is measure-zero relative to $\nu$'') --- strict dominance, not a preorder, the base case being irreflexive everywhere --- with: transitivity; \emph{asymmetry} ($\mu \ll \nu$ forbids $\nu \ll \mu$; irreflexivity follows), stipulated because without it the degenerate model with all nonzero weights mutually negligible satisfies every other axiom and collapses the semantics; $0 \ll \nu$ for every $\nu \neq 0$; relative downward closure ($\mu' \leq \mu \ll \nu \Rightarrow \mu' \ll \nu$); second-argument monotonicity ($\mu \ll \nu$ and $\nu \leq \nu'$ imply $\mu \ll \nu'$), so that negligibility relative to any weight of a combination implies negligibility relative to its total; relative additive closure ($\mu \ll \nu$ and $\mu' \ll \nu$ imply $\mu + \mu' \ll \nu$, else two invisible weights could sum to a visible one through the same-Aethus clause); scale-covariance for \emph{non-zero-divisor} $\lambda$ ($\lambda \neq 0$ sufficing at the base case) --- $\mu \ll \nu \Rightarrow \lambda\mu \ll \lambda\nu$, the restriction required since $\lambda = 0$ contradicts the base case outright ($0 \ll 0$) and a zero-divisor can collapse the two sides' levels (a two-coordinate product semiring witnesses it) --- which is the multiplicative-absorption question answered rather than left open, and the relational weight principle's second formal face, beside the cancellation identity's scale-invariance law: one governs weights under common rescaling of the combination, the other visibility under the same; and the base case $\mu \ll \nu$ iff $\mu = 0$ ($\nu \neq 0$), recovering the absolute class $\{0\}$ that Lemma~\ref{lem:basecase} consumes at $\mathbb{R}_{\geq 0}$ and making visibility context-independent there. \emph{Visibility} is then fixed by one rule across contexts: a weight $\mu$ is \emph{invisible in a context} $C$ --- a weight list, whether a combination's own or an anchoring combination's --- iff $\mu \ll w$ for some weight $w$ of $C$, and visible otherwise. On a combination's own weights this coincides with the pairwise-other reading by irreflexivity ($\mu \ll \mu$ being impossible, ``other'' does no work), it covers the anchored judgments of §\ref{sec:reduced} directly --- the anchoring context fixed to the \emph{normal-form coefficient multiset} (well-defined on congruence classes by the coefficient fact of §\ref{sec:placement2}; different representations present different lists, $aC + bC$ against $(a{+}b)C$, and the stipulation is what fixes one), relative additive closure making the stipulation harmless rather than arbitrary, since merging invisible weights cannot manufacture visibility --- and a single-term combination's weight is visible outright, the vacuous-quantifier reading excluded explicitly with the corpus's vacuity history in view; by second-argument monotonicity, invisibility against any weight of $C$ implies negligibility relative to $C$'s total wherever the total is formed --- a one-way implication. Definition~\ref{def:child-proper} compares a coefficient with the remainder's \emph{total}; the two tests coincide exactly when the measure type is \emph{sum-splitting}, $\mu \ll \sum_i w_i$ implying $\mu \ll w_i$ for some $i$, true at $\mathbb{R}_{\ge 0}$ and in graded types, adopted as a restriction on the measure class wherever the two tests must agree; otherwise the definition's total test governs proper childhood and the list test is a sufficient condition for invisibility. The unified rule's verification travels with the $\ll$ apparatus.

\textit{What the generalization is, and is not.} The character of this generalization deserves explicit statement: it is not that an Aethus \emph{gets} a measure --- rather, an accompanying object carrying a measure attaches to the Aethus per whatever construction one selects, with the original real-valued weights remaining the base case. The purpose of going beyond the reals at all is to retain gradations the reals collapse: zero weight, null-set presence, null-set presence within a null set, and so on through whatever iterated stratification a given model demands, in the spirit of non-Archimedean probability \cite{benci2013nonarchimedean} --- the negligibility relation is exactly this stratification formalized, ``null-set presence within a null set'' reading $\mu \ll \nu \ll \rho$ in chain. This is particularly suited to the self-similarity of invalidity under the consistency postulate --- the very property by which the inheritance paradox of the recursive checking Aethus (§\ref{sec:reduced}) is resolved --- and it leaves the reduced form natively prepared for fractal-dimension-style weight structure should an application demand it.

\textit{The structure on $M$, customizable; the foundational theory.} The precise algebraic structure on $M$ is left deliberately customizable, permitting ``sub-theories'' that impose additional axioms on the measure structure for specific purposes; and, like any theory of the transfinite beyond $\aleph_0$, none of this stratification is load-bearing for the foundational results --- it is capacity held in reserve. The foundational theory consists of the three postulates together with the label-to-realization convention; any additional axioms or specializations on $M$ constitute their own sub-theories.

The elements of $T$ are all \emph{finite} addition-based formal combinations (the infinite/limit extension is flagged below at its one point of use; the \emph{additive identity} $0_T$ is the empty combination) of unweighted Aethae in $U$ with weights from $M$. Every such additive combination has a representation --- each term of $M$ is not measuring something intrinsic about its attached element of $U$ but is externally imposed, a separate structural degree of freedom --- and the sum structure of $T$ is bijective to sets of (element of $U$, term of $M$) tuples, as opposed to single tuples directly, since arbitrary sums in $T$ are not reducible to rescalings of single elements of $U$. Multiplication on $T$ is the bilinear extension of $U$'s --- $(\sum_i a_i A_i) \cdot (\sum_j b_j B_j) := \sum_{i,j} a_i b_j \,(A_i \cdot B_j)$, the products evaluated in $U$ --- with its descent to the quotient below discharged in four lines in §\ref{sec:placement2}.

\begin{prop}[A coefficient system satisfying every visibility axiom]\label{prop:polymodel}
Let $M = \mathbb{R}_{\ge 0}[X]$, the polynomials with nonnegative real coefficients, with $\deg 0 = -\infty$, and set $\mu \ll \nu$ iff $\deg\mu < \deg\nu$. Then $M$ is a commutative, zerosumfree, additively cancellative semiring with unit whose algebraic order is the coefficientwise order, and $\ll$ satisfies together every axiom of the visibility list: transitivity, asymmetry, $0 \ll \nu$ for $\nu \ne 0$, relative downward closure, second-argument monotonicity, relative additive closure, scale-covariance for every nonzero $\lambda$, none being a zero-divisor, sum-splitting, and (M-comp).
\end{prop}
\begin{proof}
Nonnegative coefficients cannot cancel, so $M$ is zerosumfree, and coefficientwise subtraction makes it additively cancellative with the coefficientwise order as its algebraic order. For nonnegative coefficients $\deg(\mu\nu) = \deg\mu + \deg\nu$ and $\deg(\mu + \nu) = \max(\deg\mu, \deg\nu)$ hold exactly. Transitivity and asymmetry are those of $<$ on degrees, and $\deg 0 = -\infty$ gives $0 \ll \nu$. A coefficientwise $\mu' \le \mu$ has $\deg\mu' \le \deg\mu$, which gives relative downward closure, and likewise second-argument monotonicity. The degree of a sum being the maximum gives relative additive closure and sum-splitting, $\mu \ll \sum_i w_i$ meaning $\deg\mu < \max_i \deg w_i$. The degree of a product being the sum gives scale-covariance for every $\lambda \ne 0$. For (M-comp), $b + c = 1$ forces $b$ and $c$ to be constants, since nonnegative coefficients cannot cancel, so its premises and conclusion are read at degree zero, where they hold as at the base case.
\end{proof}

\noindent The proposition should be credited with exactly what it shows. It exhibits a coherent coefficient system with nontrivial relative negligibility, and (M-comp) on the normalized decompositions that exist in it; but those are constant-coefficient, and a nontrivially graded combination such as $1 + X$ cannot be normalized to total weight one within $\mathbb{R}_{\ge 0}[X]$, division by $1 + X$ being unavailable. The limitation belongs to the coefficient system, and a modest extension removes it. Let $M' = \{P/Q : P, Q \in \mathbb{R}_{\ge 0}[X],\ Q \neq 0\}$, with equality as rational functions, $\deg(P/Q) = \deg P - \deg Q$, and $\mu \ll \nu$ iff $\deg\mu < \deg\nu$. Sums and products of such quotients keep nonnegative coefficients, so no leading terms cancel, and $\deg(\mu + \nu) = \max(\deg\mu, \deg\nu)$ and $\deg(\mu\nu) = \deg\mu + \deg\nu$ hold as in the polynomial case; the degree identities of the proof carry over, and with them the normalized-degree argument for (M-comp). Graded normalized weights now exist: $b = 1/(1+X)$ and $c = X/(1+X)$ satisfy $b + c = 1$ and $b \ll c$. The extension strengthens the coefficient-system example, and it supplies nontrivially graded normalized coefficient chains. With $b = 1/(1+X)$ and $c = X/(1+X)$ as above, and $d = X/(1+X)$ and $e = 1/(1+X)$ splitting the component of weight $c$ in turn, $c$ and $d$ are visible against their respective remainders, the composed coefficient $cd = X^{2}/(1+X)^{2}$ has degree $0$, and the combined remainder $ce + b = (1+2X)/(1+X)^{2}$ has degree $-1$, so the composed component remains visible, the total being $cd + ce + b = 1$. Their realization by a coherent system of carrier splittings, and the extension of the reduced-form results, remain separate obligations, as does the splitting representation problem of Proposition~\ref{prop:splitting-obstruction}.

\noindent\textit{A negligible component, worked.} In this coefficient system let $V$ be a valid carrier and $I$ an invalid one, and let $x = 1 \cdot V + X \cdot I$. $V$ is a component of $x$ with coefficient $1$ against the remainder's weight $X$, and $1 \ll X$, so $V$ is not a proper child of $x$: valid though it is, it supplies no validity floor. $I$ is the only proper child of $x$, since $X \not\ll 1$, and it is invalid, so $x$ has no nonzero reduced form and is invalid, while $V$ stays represented in $x$ with its weight and is not invalidated. The example needs exactly the assumption that marks it, a coefficient system in which a nonzero weight can be negligible relative to another. $\mathbb{R}_{\ge 0}$ is not one: at the main domain every positively weighted valid carrier is a proper child and supplies a floor.

\subsection{The factorization identity and the quotient: importing the graph of $U$ into $T$}

We now give the mechanism by which the graph-theoretic content of the intrinsic layer is re-imported into the additive structure of the extrinsic layer. It splits into two parts of different status --- a derived \emph{identity} (basic factoring, no imposition at all) and the quotient proper, which is smaller than a first formulation suggests.

Given unweighted Aethae $A, B \in U$ with weights $a, b \in M$, forming $aA + bB \in T$, let $C \in U$ be the last common ancestor of $A$ and $B$ in the parent--child graph, which by construction is the informational union $C = A \cup B$. By the parent--child structure of $U$ under its multiplication, there exist $C_A, C_B \in U$ with
\begin{equation}
A = C_A \cdot C, \qquad B = C_B \cdot C,
\end{equation}
the $C_A, C_B$ encoding the additional specifying information distinguishing $A$ and $B$ from their shared parent. Solutions to $y \cdot C = A$ are not unique as Aethae --- $y = A$ itself always solves it, and mixed solutions exist besides --- and the hierarchy of what depends on which should be fixed before the semantics is chosen: the factorization identity itself is a theorem generally, at any factor solutions; the \emph{canonical} factorization is a theorem in the finite distributive case (finiteness alone does not suffice); and the general canonical residue is open, an optional canonicalization the headline result does not consume. The decomposition, where canonical, is fixed by \emph{coarsest-solution semantics}:

\begin{defin}[Coarsest-solution decomposition]\label{def:coarsest}
$C_A$ is the \emph{coarsest} $y \in U$ --- the greatest in the declared orientation, where the parent direction is the larger --- satisfying $y \cdot C = A$: concretely, the informational union of all solutions, the co-Heyting--style difference of $A$ by $C$ read at the correct end of the order. (The opposite, ``least'' orientation would be backwards on this paper's convention: every solution is refined by $A$, so the refinement-least solution is $A$ itself --- precisely the degenerate pick the semantics exists to exclude.) Existence: under Hypothesis (D), the informational union of two solutions is a solution --- $(y + y')\cdot C = y\cdot C + y'\cdot C = A$ --- so the solution set is join-closed and possesses a coarsest element in the finite case; existence in general is a booked lemma, not assumed. Coarsest-ness is what the construction below consumes; nothing here claims solutions are unique as Aethae.\footnote{An aside for readers in rich logics: under full classical consequence, the canonical difference in the counterexample below is $\mathrm{Cl}\{p \rightarrow q\}$ rather than $\mathrm{Cl}\{q\}$, since $q \vdash p \rightarrow q$ makes the former strictly coarser --- the co-Heyting difference in such logics is the implication-residue.}
\end{defin}

\noindent \emph{Non-uniqueness.} Solutions are not unique, and the counterexample is worth having: with $A = \mathrm{Cl}\{p, q\}$ and $C = \mathrm{Cl}\{p\}$, the Aethae $\mathrm{Cl}\{q\}$, $A$ itself, and $\mathrm{Cl}\{p \rightarrow q\}$ are three inequivalent solutions; the erroneous step --- from ``equivalent modulo the shared $C$-content'' to equality as Aethae --- was exactly the gap. The coarsest-solution semantics above is the repair, and it is also what makes the Case-2 trichotomy below true as stated.

Earlier formulations imposed two cases here, partitioned by whether the decomposition yields $C_A \neq C_B$ or $C_A = C_B$; the first is now the following identity, the second the quotient's sole clause.

\begin{prop}[Last-common-ancestor factorization identity]\label{prop:factor}
For any $A, B$ as above,
\begin{equation}
aA + bB \;=\; C \cdot (a\, C_A + b\, C_B), \label{eq:main}
\end{equation}
already at the level of formal combinations: bilinear multiplication evaluates the right side to $a\,(C_A \cdot C) + b\,(C_B \cdot C) = aA + bB$, the products landing in $U$ by the factorization. Nothing is removed and nothing is prevented; the last common ancestor is exhibited once, extracted as a common factor. Two scope notes: at $A = B$ the identity reads $aC + bC = C\,(a\Null + b\Null)$ and holds the same way, the earlier $A \neq B$ restriction being a fossil of the exclusive-cases framing; and the identity is decomposition-independent --- any solutions of $y \cdot C = A$ serve --- coarsest-solution semantics canonicalizing the \emph{exhibited} factors rather than enabling the identity.
\end{prop}

\noindent The quotient proper is therefore generated by a single collapse clause together with the deletion convention below --- the \textbf{same-Aethus clause} (the fuller development's Case 2):
\begin{equation}
aC + bC \;\equiv\; (a+b)\, C. \label{eq:degenerate}
\end{equation}

The identity's scope and the clause's scope meet cleanly. Under $C = A \cup B$, the same-Aethus situation $C_A = C_B$ occurs only when $C_A$ and $C_B$ are each the blank Aethus --- equivalently only when $A = B = C$: for $C_A = C_B$ forces $A = C_A \cdot C = C_B \cdot C = B$, so $C = A \cup A = A$, and then $A = C_A \cdot A$ forces $C_A = \Null$. The three conditions ``$C_A = C_B$,'' ``$A = B$,'' and ``both equal $\Null$'' are therefore equivalent here --- noting that the step ``$A = C_A \cdot A$ forces $C_A = \Null$'' is exactly where coarsest-solution semantics earns its keep: $C_A = A$ also solves the equation, and $\Null$ --- solving it and coarsest overall --- is what the semantics selects, while for $A \neq C$ the blank Aethus fails the equation ($\Null \cdot C = C \neq A$) and is correctly excluded.

The same-Aethus clause is the requirement that repeated weightings of the same unweighted Aethus combine additively --- $a$ and $b$ applied to the same $C$ yield a single weight $(a+b)$ on $C$ rather than a standing formal sum --- and it is notably \emph{not derivable} from the identity, which is \emph{uninformative} at $C_A = C_B$; substituting $C_A = C_B = \Null$ into the identity's right side returns $C\,(a\Null + b\Null) = aC + bC$, the expression to itself. Imposing the same-Aethus clause together with a \textbf{zero-term deletion convention} is what causes the formal $M$-combinations to actually constitute the commutative semiring $T$ --- the quotient is exactly same-Aethus additive collapse plus zero-term deletion, the factorization content living upstream as an identity. The deletion convention: terms of weight zero are deletable, needed because the collapse clause alone never sheds a zero term ($0 \cdot C + \mu A$ with $A \neq C$ hits the irreducible fixed point of the identity's recursion below); the ring-collapse display of §\ref{sec:placement} uses this convention to finish at $0_T$, and an ``equivalently $0 \cdot C$'' gloss follows only with this convention.

An illustrative sub-case of the identity: one of $A, B$ a parent of the other, say $aA + cC$ with $A = C_A \cdot C$. The ancestor of $A$ and $C$ is $C$ itself, the second term's decomposition trivially $C = \Null \cdot C$, and the identity gives
\[
aA + cC \;=\; C \cdot (a\, C_A + c\, \Null),
\]
which is not further reducible: applying the identity recursively to $a C_A + c\,\Null$ yields ancestor $\Null$ and returns the same expression as a fixed point. This is structurally correct --- a weighted disagreeing superposition between a parent and one of its children genuinely cannot be flattened further, and the relation correctly leaves it in canonical form.

The identity and the clause jointly re-import the graph structure of $U$ into the algebra of $T$: the last common ancestor appears explicitly in the identity as the factored-out shared parent, so formal weighted sums --- freely combinable, as §\ref{sec:placement2}'s free-construction identification says --- additionally always admit the factored form exhibiting their shared ancestry. It also settles a point of potential confusion: the addition on $T$, written $+$, is structurally analogous to the union $\cup$ on $U$ --- both combine elements via the parent--child graph --- but the two are \emph{not the same operation}. The union collapses a pair of unweighted Aethae to a single unweighted Aethus; the addition preserves weight information, every sum admitting the ancestor-factored form while remaining a free formal sum. The overlap is one of role, not identity.

\subsection{What $T$ is: the monoid semiring $M[U]$, and the congruence discharged}\label{sec:placement2}

Under the shrunk quotient, $T$'s elements are exactly the finitely supported functions $U \to M$ --- collect like terms, delete zeros --- and the bilinear product of §\ref{sec:extrinsic} is their convolution. That is: \textbf{$T$ is the monoid semiring $M[U]$}, the classical semigroup-semiring construction \cite{golan1999semirings} --- consuming, note, only $U$'s \emph{multiplicative monoid}: $U$'s addition enters solely through the factorization identity and the superstructure. Two former obligations become citations at once: normal forms and coefficient well-definedness are the standard facts about $M[U]$ (coefficients need no cancellativity; only residue uniqueness does), and the paper's bespoke content is thereby correctly located --- not in the carrier algebra, which is classical, but in the order and pruning superstructure built \emph{on} $M[U]$ from §\ref{sec:reduced} onward.

That the quotient is a \emph{congruence} for the multiplication discharges outright, in four lines, from $M$'s semiring stipulation: for $x = aC + bC + r \equiv (a+b)C + r$ and any $y = \sum_j \nu_j B_j$, termwise
\[
(a\nu_j)\,(C \cdot B_j) + (b\nu_j)\,(C \cdot B_j) \;\equiv\; ((a+b)\nu_j)\,(C \cdot B_j)
\]
by $M$-distributivity and the same-Aethus clause, while deleted zero-weight terms stay deleted by $0 \cdot \nu = 0$. The shrunk quotient never touches $U$'s lattice structure, so no appeal to Hypothesis (D) occurs anywhere in the check; (D)'s footprint in this paper is thereby a single consumer, coarsest-solution existence. (The four lines above are the actual size of the fact: the check neither needs Hypothesis (D) nor passes ``by construction.'')

What was fundamentally important in the earlier framing --- a mechanism ``forcing'' graph structure on sums that would otherwise combine freely --- is retired with the imposed-clause era: $M[U]$ \emph{is} the free construction, and the graph re-enters not by force but through the automatic factorization identity and the definitions built on top of it, the consistency postulate depending on exactly that superstructure.

One further classical identification completes the relocation, and it discharges a hypothesis the complement cluster of \S\ref{sec:placement} would otherwise consume silently. The invalid carriers form a \emph{monoid ideal} of the multiplicative monoid of $U$ by construction: with $I$ the ideal generated by the doomed set $D$ of §\ref{sec:fp}, $A \in I$ and $B \in U$ give $A \cdot B \in I$ outright, and $I = D$ by Proposition~\ref{prop:doom-tag}: with each attribute's regime stated as content, doom is closure's tag, and a tag propagates through every product, so the contraction below kills exactly the doomed carriers. It is a monoid ideal and not a semiring ideal: additive closure is not claimed and does not hold in general, since the informational union of two invalid carriers is their last common ancestor, which may be an untagged Aethus above both. That is enough. For monoid semirings the congruence generated by sending a monoid ideal's elements to zero is exactly \emph{agreement on the coefficients outside $I$}, and the quotient is the contracted monoid semiring $M[U; I]$; that is the theorem the siblings cite, and it should be cited by that name. Both inclusions are short: agreement-outside-$I$ is the kernel of coefficient restriction, a homomorphism precisely because $I$ is an ideal (a cross term with an invalid factor lands on an invalid carrier, never on a valid one), and it contains the generators; conversely the scalar chain of the congruence check above, applied termwise, places any $I$-supported difference of representations inside the generated congruence. The quotient $\bar{T}$ is therefore the classical \emph{contracted} monoid semiring $M[U; I]$ --- coefficient functions on $(U \setminus I) \cup \{0\}$ --- and four dividends fall out at once. Properness at every nontrivial $M$ is immediate, the restriction map being onto with kernel exactly the $I$-supported elements. The zero class is characterized transparently --- $x \equiv 0_T$ iff $x$ is supported on invalid carriers --- which turns the base-case forward direction of the consistency-postulate harmonization into a one-liner and exhibits the exotic residual's exact shape, congruent not to $0_T$ but to its measure-zero valid residue. \emph{Additive cancellativity of $\bar{T}$} is inherited pointwise from $M$'s stipulated cancellativity on the surviving coordinates --- the hypothesis the two-sided theorem of \S\ref{sec:placement} consumes, named there now. And the novelty scoping survives intact: the quotient, like the carrier, is a known object, and the bespoke content stays exactly where the placement section says it is.

\section{Proper children, proper Aethae, and the reduced form}
\label{sec:reduced}

\begin{defin}[Weighted proper child]\label{def:child}
An unweighted Aethus $a \in U$ is a \textbf{proper child} with weight $\omega$ of $x \in T$, written $a \lhd_\omega x$ (or $a \lhd x$), a relation of a carrier to a combination and kept visibly distinct from proper childhood $B \le A$ between carriers, the coefficient $\mu$ against the remainder's weight, read as $\mu/\lvert x \rvert$ where the measure type admits division, iff $x \ne 0_T$ and $x = \mu a + b$ with $\mu = w \ne 0$, $w$ being $a$'s coefficient in the total $x$, the supremum of the coefficients $a$ can carry over admissible decompositions, read on a chosen representation as $a$'s normal-form coefficient (Definition~\ref{def:child-proper}), and $\mu$ not measure-zero relative to the remainder's weight, $\mu \not\ll \lvert b \rvert$ in the strict relative-negligibility relation of §\ref{sec:extrinsic}, the comparison Definition~\ref{def:child-proper} makes governing; the remainder $b$ carries the rest of $x$ and can be added back. The coefficient is the maximum on the chosen representation, the weight its ratio where formed, an Aethus its own proper child at weight one. The lower condition, relative visibility against the remainder, is the load-bearing part; ``$a$'s coefficient in the total'' is well defined on congruence classes by the coordinatewise fact of $M[U]$, §\ref{sec:placement2}. Per-term \emph{droppability}, any subset of a combination's terms zeroed so long as one remains, lives in the weighted-parent relation and the nonzero restriction on $S(x)$.


\end{defin}

\begin{defin}[Proper Aethus]\label{def:proper}
$a \in T$ is a \textbf{proper Aethus} iff no $x \in U$ labeled invalid by the criteria of $T$ is a proper child of $a$ --- which, proper childhood now natively carrying weight visible against the remainder under Definition~\ref{def:child}, says exactly: no invalid Aethus appears in $a$'s combination at weight that is not measure-zero relative to the rest of the combination --- the proper-Aethus test standalone reading visibility in $a$'s \emph{own} context, consistently; only its invocation inside $S(x)$ anchors to $x$'s.

\textit{Invalid Aethae are indexed; uniqueness across constructions.} Invalid Aethae here need not be objective; they are indexed by $T$, maintaining independence from $\Hyp$ itself. For uniqueness to be preserved across different constructions of invalid Aethae --- noting that on this paper's own parent direction $\Hyp$, the maximal-information contradictory Aethus, is a \emph{child} of all Aethae rather than a parent, everything being a parent of the all-statements Aethus; and noting a question the clash-grade ruling settles rather than leaves open: deductive closure is \emph{tag-blind} --- were it explosive on tagged representations, $\mathrm{Cl}\{\text{red-all-over},\ \text{green-all-over}\}$ would be the all-statements Aethus and every tagged carrier would collapse onto $\Hyp$, contradicting Definition~\ref{def:universe}'s standing-as-elements; jointness-disallowedness is carried by tags, invisible to the consequence relation, so closure never explodes and the uniqueness at stake concerns the tag-invalid Aethae as standing carriers --- the commitment that \textbf{$T$ is a commutative semiring but not a ring} stands here, with its layered justification in §\ref{sec:placement}. Measure-zero weight in place of definitive scalar zero injects flexibility for null-set-like components to be discounted --- under the comparative definition they are simply not proper children --- a postulation making valid Aethae that much more robust in fractal-like edge cases, which may well eventually be on the table given the self-similarity of the recursive checking Aethus below.

\textit{The second payoff: the reopened blank is weighted.} A second payoff, ruled by the author on August 29, 2026 and recorded in the companion at the perturbation operator's canon \cite{Aethus}: because an invalidated component is retained at null-set presence rather than deleted, its pre-invalidation weight survives beneath the discount, and a reopening that removes the invalidating statement lifts the discount---the component returns to proper-child standing carrying the weight it carried before. The blank a reopening lands on is therefore \emph{weighted}, by the tree's own retained structure, which is the initial-condition Aethus's weighting scheme itself, (the \emph{Aetherian Nexus} of the Nexic wing); the Nexus Principle's demand that every blank carry a weighting scheme is met for reopened blanks without a second scheme, and the recovery is synchronic---the retained components are in the tree now, from every parent, not only the route walked. The condition consumed is that invalidation be indexed to its invalidators, which the joint-invalidity relation supplies; the recovery lemma---that removing an entry lifts exactly the discounts its joint invalidities imposed---is the named upgrade from ruling to theorem, and its domain is the retained pre-quotient data, the carrier with its invalidators and their coefficients, which the architecture keeps, and not the validity quotient alone: for a valid $V$ and an invalid $I$, $[V + I] = [V + 2I] = [V]$ in the quotient, so a coefficient discarded there cannot be reconstructed from the class, and recovery is from the enriched state and never from the reduced form.
\end{defin}

A translation note belongs here, since the semiring constraint determines which classical instrument applies. In a ring, one would say the invalid Aethae generate an \emph{ideal} and the partition theorem of the fuller development is computation modulo that ideal; in a commutative semiring, additive inverses unavailable by the constraint just stated, the correct instrument is a \emph{congruence} \cite{golan1999semirings}: the relation generated by identifying every invalid element of $U$ with the additive identity --- the $U$-level generator set the harmonization and kernel argument below already assume, the stipulation recorded as the author's to override --- is a commutative semiring congruence on $T$ --- in the author's formulation, the \emph{removal} of invalid Aethae is just quotienting by validity equivalence: Aethae themselves are not valued by it, which is exactly why the tags can remain standing short of the quotient .

\textit{Validity equivalence, exactly.} ``Validity equivalence'' here names the \emph{generated} congruence throughout, not the two-block partition (valid with valid, invalid with invalid), which fails to be a congruence outright: for tagged $a, b, d$ with $a + d$ the untagged blank while $b + d$ is tagged, $a \sim b$ two-block yet the sums split classes. A corollary worth keeping falls out en route: invalid inheritance's contrapositive makes every parent of a valid Aethus valid, so validity is upward-closed under $+$ --- the two facts mutually consistent, the counterexample's invalid-plus-invalid-equals-valid sitting in the direction the corollary does not constrain. --- the partition theorem is precisely computation modulo it, and, at the base case $\mathbb{R}_{\ge 0}$ under the stated hypotheses, the reduced form is the canonical representation of an Aethus's congruence class.

\textit{Where probability lives.} Two roles are kept apart. Semiring homomorphisms from $T$ that respect the invalidity congruence are \emph{algebraic valuations}: natively, an admissible \emph{weight valuation} is $M$-valued, a congruence-respecting normalized semiring valuation from $T$ into $M$ itself (a semiring homomorphism supplies additivity and multiplicativity, not normalization, so the clause is genuinely extra), the framework's postulate consuming only internal comparisons in $M$; and one such homomorphism witnesses that the quotient is nontrivial. That the congruence is \emph{proper} is discharged in general --- for any nontrivial $M$ --- by evaluation at the blank: $x \mapsto x(\Null)$ is a semiring homomorphism $T \to M$ at \emph{any} $M$ (additivity and zero coordinatewise; multiplicativity the $U$-level fact that $A \cdot B = \Null$ forces $A = B = \Null$, the closure of the union of representations containing each factor's closure inside $\mathrm{Cl}(\emptyset)$ --- the return inclusion trivial, the full chain in Lemma~\ref{lem:basecase} --- independent of $M$; the unit to $1$), its kernel a congruence containing every generating pair $(1 \cdot X,\, 0_T)$ given the blank untagged, hence the generated congruence, while $1 \cdot \Null \mapsto 1 \neq 0$ exhibits non-totality whenever $M$ is nontrivial --- Lemma~\ref{lem:basecase}'s evaluation homomorphism being this argument's base-case instance. Without properness the canonical-representation reading has no content. \emph{Empirical probability} is the other role: the specified weighting, normalized on admissible alternatives, which a model supplies when it must emit real-valued frequencies. It is additive across exclusive alternatives, and that additivity does not make it multiplicative for arbitrary joint assertions: the empirical probability assignment is not, in general, a semiring homomorphism for joint assertion, since for a real-valued $h$ multiplicativity gives $h(A) = h(A \cdot A) = h(A)^2$, forcing $h(A) \in \{0, 1\}$ on every idempotent carrier (a general coefficient semiring can have other idempotents). A second requirement is independent of the first: the probability of realized content must respect the validity identification. If a model admits $B + I$ as a cell beneath $A$ and not $B + J$, with $A$ and $B$ valid and $I$ and $J$ invalid, the two still represent one class, $[B + I] = [B + J] = [B]$; raw accessibility distinguishes representatives of the same surviving content, so the probability of realized content is evaluated on classes, a raw-level model of the unfolding inducing it by collecting the histories that realize each class \cite{Aethus}. Which weighting it is --- the valuation question, \emph{which} valuation centric unfolding follows, and why --- is answered by the weighting scheme every blank carries, the Nexus Principle's demand, and it is one question for every attribute: by the long paper's Randomness Universality Principle, a quantum outcome is no different in kind from any other centric-unfolding update, only easier to parse as blank \cite{Aethus}. For quantum attributes the scheme is the Quantum Nexus kernel, the one flagged kernel axiom of the physics-face companion \cite{AethusQuant}, whose weights are the Born weights, so that the Born rule is that kernel's corollary and not a separate question. Centric unfolding prescribes the eventual outcome probabilities of a specified resolution, its weights in the validity quotient, while leaving the finite-step process underdetermined within its admissible family \cite{Aethus}. Unnormalized mass is additive across exclusive alternatives, and normalized probability is linear in normalized mixture coefficients with the relevant denominator fixed; it is not in general a linear function of arbitrary raw coefficients. The weights enter through the weighting scheme on terms, through the kernel for quantum attributes, and additivity holds across record-separated terms, not across interfering alternatives inside one agreeing term.

Let $x \in T$ be a \textbf{weighted parent} of $y \in T$ with \textbf{parent residue} $z$ iff $x = y + z$ for some $z \in T$ with all composing measures of $y$ and $z$ non-negative --- a rider automatic under the stipulated algebraic order (every coefficient is non-negative definitionally), retained for readability across general cones; the substantive guard is zerosumfreeness (§\ref{sec:placement}).

\begin{defin}[Reduced form]\label{def:reduced}
Partially order the set $S(x)$ of \emph{nonzero} proper Aethae in $T$ taking $x$ as weighted parent --- ``nonzero'' read \emph{modulo the visibility structure, anchored to $x$'s context}: a candidate counts as nonzero for membership in $S(x)$ only if some carrier of it is \emph{visible relative to $x$'s weights} --- own-context visibility is refuted by an infinitesimal counterexample (take $\varepsilon \ll 1$ and $x = \varepsilon V + 1 \cdot I$, $V$ valid, $I$ invalid: in $\varepsilon V$'s own context $\varepsilon$ is the sole weight, hence visible, and the invisible carrier would rescue $x$'s validity, exactly the aid the semantics forbids; anchored to $x$, $S(x)$ is empty and $x$ invalid, as required) --- with the semantics stated exactly: measure-zero-weight carriers are not thereby invalidated (no tag flips), they simply fail to be \emph{visible} to the postulate's recursive condition over $T$, unable to come to the aid of preventing invalidity when it comes down to it --- the restriction to nonzero Aethae remaining load-bearing after the comparative clarification, for a simpler reason than before: $0_T$, having no visible components at all, is vacuously proper under Definition~\ref{def:proper} and sits in every $S(x)$ via $x = 0_T + x$; unrestricted, a reduced form would always exist and the consistency postulate never fire --- by the weighted-parent relation itself, the parent direction the ``larger'' direction (consistent with probability events, where a union of events is the larger event). The \textbf{reduced form} $\bar{x}$ is the \emph{maximum} of $S(x)$ under this ordering when the maximum exists, with $x$ declared \emph{without} a reduced form when $S(x)$ is empty; that a maximum (not merely a supremum) exists whenever $S(x)$ is nonempty is discharged at the base case by Lemma~\ref{lem:basecase} (the valid-support truncation), the exotic-measure residual booked, since a canonical representation requires attainment.

\textit{Intuitively.} Intuitively, the reduced form is the proper Aethus closest to $x$ in the parent direction --- formed by pruning $x$ as little as possible while still being proper.

\textit{The consistency postulate, restated.} The \textbf{consistency postulate} (the third postulate of the companions) then states: $x$ without a reduced form --- equivalently, $x$ whose pruning bottoms out at the additive identity --- is labeled invalid by the criteria of $T$. This form harmonizes with the congruence picture of the translation-note paragraph above, with the two directions carrying different statuses, kept distinct: ``bottoms out at $0_T$'' is the \emph{definitional} gloss of $S(x)$-emptiness, while the forward direction --- invalid $\Rightarrow$ congruent to $0_T$, the $T$-level identity coincident with $\Hyp$'s class --- is a \emph{theorem}, established at the base case: by Lemma~\ref{lem:basecase} every positively weighted carrier of such $x$ is invalid, and the scalar chain $\mu X = (\mu \Null) \cdot (1\,X) \equiv (\mu \Null) \cdot 0_T = 0_T$, applied termwise and summed, lands $x$ in the zero class --- at exotic measure types $S(x)$ can conceivably empty through measure-zero structure without every carrier being invalid, so the forward direction carries the \emph{same} exotic residual as the converse, booked symmetrically. The converse --- nothing valid swept to the zero class --- is the properness fact, likewise discharged at the base case by Lemma~\ref{lem:basecase}; granted the residuals, at the base case, the reduced form of a valid Aethus is its nonzero canonical representation, the anchored criterion and the quotient parting for generalized weights.
\end{defin}

\paragraph{Sharpenings: one road not taken, one adopted, one open.} First: under the \emph{literal} reading of Definition~\ref{def:reduced}'s ``nonzero'' ($\neq 0_T$, as against modulo the measure-zero class --- the choice made below), the valid-support truncation $x_V$ is always proper and $x = x_V + x_I$ always witnesses parenthood via Lemma~\ref{lem:pointwise} at general $M$, so $S(x)$ is empty iff every positively weighted carrier is invalid \emph{at all measure types}; the scalar chain then makes the forward direction a general theorem and the zero-class residual vanishes, the surviving residual on that side being maximum-existence alone (a proper $y$ may carry measure-zero weight on an invalid carrier, escaping the pointwise bound). The reading adopted is the \emph{modulo} one, with invisibility semantics: relatively measure-zero carriers are not invalidated, they fail to be visible to the recursion, the class's axioms living with the measure type in §\ref{sec:extrinsic} in their relative form. One later generalization, now \emph{anchored} but with re-verification open: under these axioms $S(x)$ is empty iff every carrier visibly weighted in $x$'s context is invalid, the own-context form being refuted by an infinitesimal counterexample, and the zero-class residual becomes exactly characterizable, a postulate-invalid $x$ being congruent not to $0_T$ but to its measure-zero valid residue, the obstruction the symmetric booking gestures at. So that booking stands, the literal branch retained here as the record of the road not taken.

\textit{Second: congruence-properness at all nontrivial $M$.} Second: congruence-properness discharges at all \emph{nontrivial} $M$, not only the base case --- the argument now \emph{promoted} and displayed at the translation note above, independently re-verified, rather than repeated here; what is not duplicated there, and still does real work, is the dependency accounting: the discharge consumes only the generator set and the blank's untagged validity, and the $\mathbb{R}_{\geq 0}$ codomain is needed only for the admissible-probability reading --- and the position taken is that $M$-valued weight is framework-native, the third postulate consuming only internal relative comparisons in $M$; the existence of an $M \to \mathbb{R}_{\geq 0}$ homomorphism is accordingly not a framework demand but an empirical-interface question, arising exactly when a model must emit real-valued frequencies, and it travels with the valuation question, which for quantum attributes the physics-face companion's kernel answers \cite{AethusQuant}. The dependencies, disentangled: the properness sharpening consumes the generator-set stipulation and the blank's validity (were $\Null$ invalid, it would be a generator and $1 \cdot \Null \mapsto 1 \neq 0$ would break kernel-contains-generators); the admissible-probability reading additionally the exotic-$M$ homomorphism question; the emptiness sharpening the modulo pin.

\paragraph{The worked recursive example.} The canonical instance is the fuller development's recursive checking Aethus\footnote{``Oliver's Aethus'' in the fuller development \cite{Aethus, AethusBrief}: an observer at a two-slit screen who at each checking round either checks the which-slit detector or defers, recursively.} $O_1$, constructed through the self-similarity condition
\[
O_i = (1 - x_i)\, O_{i+1} + x_i\, (\alpha_i A_i + \beta_i B_i),
\]
with $x_i$ the weight of checking exactly at round $i$, and $A_i, B_i$ the left- and right-slit events. Because $O_1$ is explicitly constructed with the both-slits agreeing superposition as a parent, each of $A_i, B_i$ is jointness-disallowed against that parent --- a tagged clash, per the ruling below.

\textit{The verdict, by weights and arity.} The verdict, however, depends explicitly on the weights and on the combination arity, and an unconditional reading would be wrong at the edges. Throughout, Definition~\ref{def:proper}'s test reads the combination's components as the \emph{joint} Aethae (parent $\cdot\, A_i$), labeled invalid by the criteria of $T$, since bare slit-events are not invalid.

\textit{The grade of the invalidity.} The grade is settled: \emph{all} invalidity in this construction is label-grade --- the intrinsic jointness-disallowed combinations are stipulated tags on standing elements, and the postulate-propagated tags stand likewise; Definition~\ref{def:universe}'s multiplication keeps every carrier distinct, and the example's joint Aethae are standing tagged elements. At $x_i \equiv 0$ these carry weight zero, fail proper childhood's native strictness, and never obstruct properness: the Aethus is \emph{valid}, an observer who never checks keeping the superposition; and if infinite combinations are admitted, the never-check tail carries weight $\prod_i (1 - x_i)$, and whenever that product is positive \emph{and visible in $O_1$'s context} the tail is a candidate proper parent and validity again survives --- at the base case, positivity suffices; at exotic $M$, a positive but invisible tail ($\prod_i (1 - x_i) \ll$ the committed checking weights) fails anchored-nonzero and rescues nothing, the anchored semantics extending the invalidity verdict to positive-but-negligible tails (the fuller statement of this dividend is reserved by the author for the fuller development). The invalidity deduction holds when the checking is genuinely committed --- $x_i$ not all zero, with the tail weight vanishing ($\prod_i (1 - x_i) = 0$, e.g.\ $\sum_i x_i = \infty$) or the combination truncated finitely so no tail exists: then no iteration $O_j$ lacks a positively weighted invalid child, the pruning bottoms out at the additive identity --- $S(O_1)$ containing no nonzero proper Aethus --- and $O_1$ is deduced invalid via the postulate --- constructed as such by asserting exactly this recursive structure for $O_1 \in T$, this example being the one flagged point of use of the infinite extension.

\textit{The self-similarity motivation.} This also exhibits the central self-similarity motivation of the postulate: to take $O_1$ as invalid although its invalidity is weaker than an explicitly asserted contradiction, since only the $A_i, B_i$ are ever explicitly deduced contradictions within the weighted layer. (``Explicitly deduced'' here means explicitly deduced \emph{tags}: the clash of stated value against maintained agreeing superposition is label-grade like everything else in the construction, the flag that formerly stood here being resolved --- carriers stand, tags propagate, and the earlier collapse reading is retired.)

\paragraph{Label versus realization; core and branches.} Within this formalism ``invalidity'' is a label assigned by the criteria of $T$, not a realized empirical property; the connection to the empirical forbidding of realization is the separately injected principle of §\ref{sec:scaffolding}. In its sharpest architectural form: invalid Aethae are perfectly consistent with the graph itself --- passively tagged, not contradictions within it --- and the consistency postulate operates over $T$, deducing further tags from the original tags. Reality then ``extracts'' at the very end: centric unfolding (the realization dynamics of the fuller development) is disallowed into tagged cases.

\textit{Backend and frontend.} The picture is backend versus frontend: $T$ in the backend remains fully well-defined no matter how the invalidity tags are placed, while reality in the frontend simply reads $T$ and prevents the agent's centric unfolding into invalid Aethae --- \emph{why} being perhaps the remaining central philosophical mystery, with the phenomenological law of noncontradiction itself \emph{derived} from this arrangement: the existing mystery taken and buried deeper into formal ontology.

\textit{Core and branches.} The three postulates plus that convention constitute the \emph{core} theory; additional axioms --- on $M$, on $U$, or as further invalidity-labeling postulates --- generate \emph{branch} theories specialized for particular purposes. Two major branches already developed in the broader framework: one extending the core with geometric apparatus aimed at deriving quantitative quantum rules as inferred consequences, and the anthropic branch extending the core to systems constrained by selection effects (yielding, for example, a derived path from a mediocrity principle to the Rare Earth hypothesis \cite{NexicFoundation}). Crucially, a branch need not be motivated by any ontological principle: one may append a validity-labeling postulate as a purely formal extension, for toy exploration or stress-testing, and this too is a branch. Recursively, even the consistency postulate itself can be viewed as one such added labeling postulate --- added precisely to correct the recursive checking setup to match empiricism --- and the choice to halt at three postulates rather than continue adding is itself a \emph{hypotheses non fingo} \cite{newton1687principia} move at the meta-level: the core deliberately refuses further invalidity axioms beyond empirical fit at the targeted level of phenomena.

This label-versus-realization distinction is also exactly why the postulate is stated as an implication, $V[A] \Rightarrow \exists B \le A,\ \forall C \le B,\ V[C]$, rather than a biconditional: its operative content is to propagate invalidity labels via its contrapositive (no reduced form $\Rightarrow$ invalid), never to derive physical realizability. Formal validity is another matter: it is the complement of generated doom, and \S\ref{sec:fixedpoint} characterizes it (Lemma~\ref{lem:fp}). What the formalism does not commit to is the further claim that a formally valid Aethus is physically realized, validity being, in that sense, the residue of no label having been assigned --- a problem-of-induction-style absence rather than a positively derived realization; asserting more would be an empirical-realization claim beyond what the formalism commits to, avoided under the same \emph{hypotheses non fingo} discipline. Correspondingly, ``contradiction'' --- the base case of the recursion --- is the exhaustive totality of clashes within the intrinsic compatibility structure of attributes, from Boolean clashes to circuit-level joint clashes; the attribute is the structural carrier of all of it, with any particular formal apparatus (such as the identifier--class--relation--set tuples of the fuller development) capturing the carriage without claiming to exhaust the abstraction, in the same sense that no single algebraic definition of a tensor quite captures the object among all its formalizations.

\section{The fixed-point lemma}\label{sec:fp}
\label{sec:fixedpoint}

The recursive reading of the consistency postulate promises a formal equivalence between the postulate as stated and a greatest-fixed-point formulation. The lemma, with branching conditions explicit:

\begin{lemma}[Fixed-point form of the consistency postulate]\label{lem:fp}
Work over the class of Aethae ordered by the \emph{declared descent} $\le$ of Definition~\ref{def:descent}, reflexive and transitive, so that the operator's quantifiers range over declared descendants at every depth. Throughout this lemma, $\le$ denotes proper childhood, the relation of Definition~\ref{def:child-proper} on $U$ under the name declared descent. It is distinct from the ambient lattice order $\sqsubseteq$ of Proposition~\ref{prop:order}, which would place the tagged all-statements carrier beneath every carrier and make the resulting doom set universal, $G(\mathrm{Base}) = U$. The weighted proper child of Definition~\ref{def:child} is the same relation read on combinations in $T$, the sense in which the two readings coincide being Proposition~\ref{prop:harmonize}. Let $\mathrm{Base}$ collect the non-recursive modes of invalidity (direct contradictions and second-postulate-contrapositive violations). Define the safety operator on sets of Aethae by
\[
F(Y) \;=\; \{\, A \notin \mathrm{Base} \;:\; \exists B \le A,\ \forall C \le B,\ C \in Y \,\},
\]
and the dual doom operator by $G(X) = \mathrm{Base} \cup \{A : \forall B \le A,\ \exists C \le B,\ C \in X\}$. Then: (i) $F$ and $G$ are monotone, so the greatest fixed point $\nu F$ and least fixed point $\mu G$ exist for arbitrary branching and are complements; (ii) the postulate as stated asserts exactly that the validity predicate $\mathcal{V}$ is a post-fixed point of $F$, $\mathcal{V} \subseteq F(\mathcal{V})$, whence $\mathcal{V} \subseteq \nu F$ always; (iii) if invalidity is generated by the base cases --- every invalid Aethus invalid either directly through $\mathrm{Base}$ or by upward propagation of the two-quantifier rule, the standing recursive reading --- then $\mathcal{V} = \nu F$; and (iv) because $\le$ is transitive, both fixed points are reached at the first stage for any branching whatever: $\mu G = G(\mathrm{Base})$ and $\nu F = F(\complement\,\mathrm{Base})$, so an Aethus is doomed iff every descendant has a base case beneath it and---under the base-generation assumption of clause (iii), without which only $\mathcal{V} \subseteq \nu F$ is proved---valid iff some descendant has a base-free cone, and the postulate's $\exists\forall$ statement is the fixed point itself. \end{lemma}

\begin{proof}
(i): if $Y \subseteq Y'$, any witness proper child $B$ of $A$ with all proper children in $Y$ has all proper children in $Y'$, so $F(Y) \subseteq F(Y')$; $G$ is identical with quantifiers exchanged. Knaster--Tarski \cite{tarski1955lattice} supplies $\nu F$ as the union of post-fixed points and $\mu G$ as the intersection of pre-fixed points, with no finiteness assumption. Complementarity is a De Morgan computation: $A \notin F(\complement X)$ iff $A \in \mathrm{Base}$ or every proper child $B$ has some proper child $C \in X$, i.e., $\complement F(\complement X) = G(X)$; taking fixed points identifies $\complement\,\nu F = \mu G$. (ii): the postulate's implication together with disqualification of the base cases says precisely $\mathcal{V} \subseteq F(\mathcal{V})$; $\nu F$ being the largest post-fixed point, $\mathcal{V} \subseteq \nu F$ follows. (iii): base-generation says the doomed set $D = \complement\,\mathcal{V}$ contains nothing beyond what $\mathrm{Base}$ and propagation force, i.e., $D$ is the least pre-fixed point of $G$, so $D = \mu G$; complementing gives $\mathcal{V} = \nu F$. (iv): let $D_1 = G(\mathrm{Base})$ and suppose $A \in G(D_1)$. If $A \in \mathrm{Base}$ then $A \in D_1$ immediately; otherwise membership in $G(D_1)$ supplies the quantified condition, that every $B \le A$ has some $C \le B$ with $C \in D_1$. If $C \in \mathrm{Base}$, $B$ has a base case beneath it; otherwise $C \in D_1 \setminus \mathrm{Base}$, so $C$'s own reflexive refinement supplies $E \le C$ in $\mathrm{Base}$, and by transitivity $E \le B$. Either way every $B \le A$ has a base case beneath it, i.e., $A \in D_1$; so $G(D_1) = D_1$, and $D_1$, the first stage after seeding with $\mathrm{Base}$---the second application of $G$ from $\emptyset$, since $G(\emptyset) = \mathrm{Base}$---is the least fixed point. Complementing through (i) gives $\nu F$. No branching hypothesis is used. (Under an immediate-child reading of $\le$ the closure ordinal would be bounded by $\omega$ under finite branching, by K\H{o}nig's lemma, and stages beyond $\omega$ would be populated under infinite branching; for the transitive relation the ordinal is one under the seeded convention $D_0 = \mathrm{Base}$, and none of that machinery is needed.) \end{proof}

\noindent\emph{The construction, in one place.} The third postulate has one definition in this paper. On $U$, an Aethus is \emph{valid} iff it lies in $\nu F$ over declared descent, clause (iii) being adopted as the definition of invalidity rather than a hypothesis about it: the tags are the base cases, doom is their propagation, and an Aethus is invalid iff it is tagged or every declared descendant of it has a tagged declared descendant. On $T$, a weighted Aethus survives iff some proper child of it, a term carried at weight not measure-zero relative to the rest of the combination, is a valid Aethus --- at the base case $\mathbb{R}_{\ge 0}$, some term of positive weight; that is the reduced-form criterion of §\ref{sec:reduced}, read through the weighted proper children of Definition~\ref{def:child}, the terms whose coefficients are not measure-zero relative to the rest. Proper childhood is read in two places, on carriers (declared descent) and on combinations (weighted terms at their full weights); Proposition~\ref{prop:harmonize} states the sense and the scope in which the two readings are one postulate.

\noindent\emph{The multiplicative ideal.} Let $D = \mu G$ be the doomed set and $I$ the multiplicative ideal of $U$ that $D$ generates, the set of products $A \cdot B$ with $A \in D$. The validity congruence of §\ref{sec:placement2} kills $I$, so the quotient is proper iff $\Null \notin I$, and $\Null = A \cdot B$ forces $A = \Null$, so properness is exactly the validity of the blank Aethus, which holds whenever some single-valued completion of it is untagged, by Proposition~\ref{prop:descent}. Doom is hereditary along declared descent by transitivity; that it is hereditary along every joint assertion, so that $I = D$, is the content of the next proposition, which rests on declared descent as rules (a)--(c) now state it, fillings declaring nothing, and on closure's splitting rule read over proper children, a carrier holding a carving each of whose settings of visible weight (its proper children) is tagged being tagged --- the filter the postulate's blindness requires, a negligible setting being unable to rescue the parent and so unable to shield it. The hypotheses the proposition consumes are these, stated together: \emph{termination}, that the model has finitely many attributes, so that successive choices of untagged settings reach a terminal descendant in finitely many steps; \emph{declared descent}, as rules (a)--(c) state it, a filling declaring nothing and an agreeing attribute declaring only the splits its access licenses; \emph{tagging}, that $\mathrm{Base}$ is the set of tagged carriers and closure's splitting rule holds; and the \emph{descent proposition}, Proposition~\ref{prop:descent}, for the validity of a blank with an untagged completion.

\begin{prop}[Doom is the tag]\label{prop:doom-tag}
In a model with finitely many attributes, a carrier is doomed iff it is tagged, and the doomed set is a multiplicative ideal: $I = D$.
\end{prop}

\begin{proof}
Every declared split is a carving's: under rule (b) directly, and under rule (c) because the access carrier records the attribute's value and so holds its carving. Tagged carriers lie in $\mathrm{Base} \subseteq D$. Conversely let $A$ be untagged. Every carving $A$ holds has an untagged setting of visible weight, since otherwise the splitting rule, read over proper children, would tag $A$; that setting is an eligible proper child, not merely a candidate component, so pass to the declared child it gives, untagged, and repeat.  Each declared split fixes an attribute, and a fixed attribute is never split again, so with finitely many attributes the chain ends, at an untagged carrier $B$ with no declared children. Its only declared descendant is itself, so $\forall C \le B,\ C \notin \mathrm{Base}$, and $A \notin D$. Thus $D$ is the tagged set. A tag is carried by content, and $A \cdot B$ holds the content of $A$, so $A \in D$ gives $A \cdot B \in D$ for every $B \in U$: $D$ is closed under multiplication by $U$, and $I = D$.
\end{proof}

\noindent The hypothesis that fillings declare nothing is doing the work. Read with an unstated attribute as an exclusive blank by default, as this section was first written, the construction admits doomed carriers that closure does not tag, $X1$ with $b$ unstated among them, whose products with an irreducibility marker are valid; $I$ then strictly extends $D$, the semiring quotient kills $I$, and the reduced form of §\ref{sec:reduced}, a map of additive monoids and not of semirings, removes exactly $D$. Letting broad agreeing states split pin by pin opens the same gap one level up, as the extended model of §\ref{sec:model} shows. With regimes as content, the reduced form, which at the base case $\mathbb{R}_{\ge 0}$ removes exactly $D$, coincides there with the semiring quotient; for generalized weights the anchored criterion can reject a combination whose quotient class survives, and a maximum reduced form can fail to exist, so the coincidence is a base-case statement.

Three remarks situate the lemma. First, hypothesis (iii) restates the framework rather than adding to it: the recursive propagation reading, base cases as leaves and the two-quantifier rule as the recursive step, is its own account of how one reduced-form application traverses descendant structure, so $\mathcal{V} = \nu F$ is the postulate's content once that reading is adopted. Second, the physics-face companion \cite{AethusQuant} carries the finitely-branching special case; this lemma is the general form, the companion's proposition its clause (iv) restricted. Third, the lemma does not by itself make the chess idiom literal: membership in $\mu G$ says every proper child has a tagged proper child, weaker than every continuation meeting a tag in finitely many steps, since a tree in which each untagged node $a_n$ has an untagged proper child $a_{n+1}$ and a tagged one $b_n$ is finitely branching, puts $a_0$ in $G(\mathrm{Base})$, and carries the infinite path $a_0, a_1, \ldots$ meeting no tag. So the fixed point's stages do not count generations of forced loss, and a game whose losing positions are exactly $\mu G$ would have to be defined and the correspondence proved. What holds, under the finite-attribute hypothesis of Proposition~\ref{prop:doom-tag}, is finite termination: each declared split fixes an attribute for good, so every chain of splits ends and doom is decided in finitely many of them.


\subsection{The weighted and unweighted formulations, harmonized}
\label{sec:harmonize}
The paper has called two things the third postulate: the $\exists\forall$ refinement-safety rule on the declared descent of Definition~\ref{def:descent}, which §\ref{sec:fp} put in fixed-point form, and the reduced-form criterion on $T$ of §\ref{sec:reduced}, which runs on additive terms at their full weights and weighted parenthood. The two are one relation, proper childhood (Definition~\ref{def:child-proper}), read once on carriers and once on combinations; the theorem owed says in what sense the two criteria are one postulate, a representation statement about the free algebra rather than a reconciliation of two primitives. The sense is bounded by two hypotheses the theorem states and the later summaries carry: the forward implication for superpositions holds when the carriers are declared descendants of the support's last common ancestor, and the converse under cofinality. It is stated at the base case $M = \mathbb{R}_{\ge 0}$, where §\ref{sec:reduced} proves that the reduced form of $x$ is the truncation of $x$ to its formally valid support when that truncation is nonzero, and $0_T$ otherwise; the exotic-measure residual is left where §\ref{sec:reduced} leaves it.

\begin{prop}[Harmonization]
\label{prop:harmonize}
Let $\iota : U \hookrightarrow T$, $A \mapsto 1 \cdot A$, be the embedding, let $D = G(\mathrm{Base})$ be formal invalidity on $U$ over declared descent, and let $x = \sum_{i} \mu_i A_i \in T$ have finite support with every $\mu_i > 0$, $L := \sum_i A_i \in U$ its support's last common ancestor. At the base case:
\begin{enumerate}
\item[(i)] \textup{(Pure Aethae, exact.)} $\iota(A)$ has a nonzero reduced form iff $A \notin D$, i.e., iff $\exists B \le A\ \forall C \le B,\ C \notin \mathrm{Base}$. The weighted criterion on an embedded pure Aethus is the unweighted criterion on the Aethus.
\item[(ii)] \textup{(Superpositions, forward.)} $x$ has a nonzero reduced form iff some carrier $A_i \notin D$, i.e., iff the $\exists\forall$ criterion holds at some carrier; and, when the carriers are declared descendants of $L$, as they are when $x$ is a disagreeing superposition of fillings of an exclusive blank of $L$, this implies $L \notin D$, since then $A_i \le L$ and a safe cone beneath $A_i$ is a safe cone beneath $L$.
\item[(iii)] \textup{(Superpositions, converse, under cofinality.)} If the carriers are \emph{cofinal} beneath $L$---every declared descendant of $L$ has a declared descendant lying beneath some $A_i$---then $L \notin D$ implies that some $A_i \notin D$, hence that $x$ has a nonzero reduced form: a safe cone beneath $L$, at $B$, contains a descendant $B' \le B$ lying beneath some $A_i$, and the cone beneath $B'$ is safe. (Requiring every descendant of $L$ itself to lie beneath some $A_i$ would be degenerate: $\le$ is reflexive, so $L$ would have to lie beneath some $A_i$, forcing $A_i = L$. Where descent terminates, cofinality is terminal coverage: every terminal descendant of $L$ lies beneath some $A_i$.) Without cofinality the converse fails: $L$ may own a safe cone none of whose members has a descendant beneath any of the $A_i$, in which case $L$ is formally valid while $x$'s reduced form is $0_T$.
\end{enumerate}
\end{prop}
\begin{proof}
(i) The weighted proper child of $1 \cdot A$ is $A$ itself at weight $1$, its supremum coefficient, every other carrier having coefficient $0$ in $1 \cdot A$ and so being no proper child of it; the reduced form is $1 \cdot A$ if $A \notin D$ and $0_T$ if $A \in D$ by the base-case truncation, and $A \notin D$ is the $\exists\forall$ condition by §\ref{sec:fp}. (ii) The truncation to valid support is nonzero iff some $\mu_i A_i$ survives, i.e., some $A_i \notin D$; the implication to $L$ is transitivity of $\le$. (iii) If $L \notin D$ then some $B \le L$ has a base-free cone; by cofinality some $B' \le B$ lies beneath some $A_i$, and the cone beneath $B'$ is base-free, so $A_i \notin D$. The failure without cofinality is the stated configuration, and it is not exotic: it is the ordinary case of a superposition that omits an alternative its parent still admits.
\end{proof}



\subsection{A finite example, worked through}\label{sec:model}

\noindent One finite example at the main domain, followed from carriers to the quotient; it is a fragment of a closed working universe, not the universe itself. Two attributes, $a$ and $b$, each with values $1$ and $0$, and nothing declared between them, so that the carvings are closure-independent; a carrier $P$ holding the single-pin carvings of both, so that its two blanks are exclusive and split (rule (b)); and parent-conditional weights for $P$'s four finest cells, $w(a1b1) = 0.2$, $w(a1b0) = 0.3$, $w(a0b1) = 0.4$ and $w(a0b0) = 0.1$.

\emph{Carriers.} The example follows a fragment: $P$; its eight cells $P_f$, one for each nonempty partial filling $f$ of $a$ and $b$; and three tagged carriers, $T_a = P_{a1} \cdot P_{a0}$ and $T_b = P_{b1} \cdot P_{b0}$, each holding rival values of one attribute, and the all-statements carrier $\Hyp$. Under the paper's non-explosive tags $T_a$ keeps its content and is not the all-statements carrier. The fragment sits inside a working universe closed under the intrinsic operations, which contains further carriers; it is not itself claimed to be closed.

\emph{Declared splittings.} The carving of $a$ gives $P \rightsquigarrow 0.5\, P_{a1} + 0.5\, P_{a0}$, the carving of $b$ gives $P \rightsquigarrow 0.6\, P_{b1} + 0.4\, P_{b0}$, and their composition gives $P \rightsquigarrow \sum_{i,j} w(a_ib_j)\, P_{a_ib_j}$. The example satisfies (S1)--(S3), the finest cells extending a cell summing to its coefficient, $0.2 + 0.3 = 0.5$ for $P_{a1}$ and $0.2 + 0.4 = 0.6$ for $P_{b1}$. It satisfies closure-independence because nothing is declared across $a$ and $b$: each partial filling closes to a distinct carrier, and content inclusion between cells is extension of fillings.

\emph{Proper children.} Every weight is positive, so every cell is visible, and since the two carvings are all the splitting $P$ declares, the proper children of $P$ are $P$ and its eight cells, the full cone and the join-semilattice $\mathcal{P}_P$ of Theorem~\ref{thm:proper-lattice}. All of their pairs have joins: $P_{a1} \vee P_{a0} = P$, two cells of one carving joining at the parent. Meets exist only for compatible fillings, $P_{a1} \wedge P_{b1} = P_{a1b1}$, their joint assertion; $P_{a1}$ and $P_{a0}$ have no common proper child. Their joint assertion $T_a$ exists in the universe, but it holds rival values of $a$, arises from no declared splitting, and is a proper child of nothing but itself (rule (a)); only in the bottom-adjoined completion do the two meet, at the auxiliary $\bot$. A child of $P$ along a carving it does not hold lies below $P$ in the intrinsic order but is no proper child of $P$.

\emph{Doom.} With finitely many attributes the doomed set is exactly the tagged set (Proposition~\ref{prop:doom-tag}): here $T_a$, $T_b$, $\Hyp$ and the other carriers of the universe that hold rival values. No cell of $P$ is doomed, each joint filling being its own safe cone. A tagged cell inside the cone would need a constraint between $a$ and $b$, and with binary attributes such a constraint breaks closure-independence, as the remarks after Theorem~\ref{thm:proper-lattice} show, so the example takes its tags from outside the cone.

\emph{Reduced form and quotient.} Let $x = 0.2\, T_a + 0.3\, P_{a1b0} + 0.4\, P_{a0b1} + 0.1\, P_{a0b0}$ in $T$. Its reduced form is the valid-support truncation $x_V = 0.3\, P_{a1b0} + 0.4\, P_{a0b1} + 0.1\, P_{a0b0}$ (Lemma~\ref{lem:basecase}), and in the validity quotient $[x] = [x_V]$, the tagged term's image being zero. In $T$ the term $0.2\, T_a$ is still present under its tag: the carrier is represented, and only its image in the quotient vanishes. In $y = 0.5\, T_a + 0.5\, \Hyp$ every positively weighted carrier is doomed, so $S(y)$ is empty, $y$ is invalid and $[y] = 0$. The accompanying script \texttt{weight\_algebra\_current\_model.py} checks closure-independence and recomputes each of these lines.

\section{Placement, and what remains open}
\label{sec:placement}

\paragraph{What is proved, and what is left open.} The paper's results divide into two classes, and the division is stated once so that no open generalization is read as part of a theorem. \emph{Proved, with no residual}: the refinement order is closed-content inclusion and a partial order (Proposition~\ref{prop:order}); joint assertion is the least common descendant and informational union the last common ancestor, the latter existing and unique for every pair (Propositions~\ref{prop:lcd}, \ref{prop:lca}); formal invalidity is the least base-generated doom set $D$ over declared descent, closed downward along it; with each attribute's regime stated as content, doom is closure's tag and $D$ is itself a monoid ideal (Proposition~\ref{prop:doom-tag}, for models with finitely many attributes), so the contracted monoid semiring quotient kills exactly the formally invalid carriers and coincides, at the base case $\mathbb{R}_{\ge 0}$, with the reduced form, which removes exactly $D$ there; the factorization identity holds at any factor solutions; the weighted algebra over a working universe is the monoid semiring $M[U_G]$; the fixed point of the consistency postulate closes at the first stage after seeding, for arbitrary branching; and the weighted and unweighted formulations agree on pure Aethae exactly, forward on superpositions, and conversely under cofinality, at the real-valued base case. \emph{Open, and consumed by none of the above}: general existence of canonical coarsest factors beyond the finite distributive case (finiteness alone does not suffice); distributivity of $U_G$ beyond the fragments whose Lindenbaum lattice is distributive; the exotic-measure residuals of relative negligibility beyond $M = \mathbb{R}_{\ge 0}$, including the anchored-emptiness characterization; the converse of the harmonization without cofinality; and the coincidence of doom and tag beyond finitely many attributes, where it needs declared descent to be well-founded. The main theorems are stated at the base case wherever a residual would otherwise enter, and the residuals are recorded here and in the sections that raise them.


Three genre placements locate the bespoke part exactly, and they are stated first because the algebra alone will read as standard to the right specialist --- which it is, and deliberately. The nearest genre is the provenance semirings of database theory \cite{green2007provenance}: weighted derivations valued in a commutative semiring, with specialization by homomorphism --- $T$'s genre precisely, the differentia being the carrier and the step, deductively closed partial Aethae under pruning rather than relational tuples under joins, with the congruence quotient doing the identification work that polynomial provenance leaves to the free object. Weighted automata and formal power series \cite{droste2009handbook} supply the second coordinate: monoid semirings of the $M[U]$ kind as behavior spaces, the differentia being the absence of any acceptance structure --- refinement, not recognition, drives the dynamics. And the fixed-point layer belongs to domain theory \cite{abramskyjung1994domain}: a refinement order with a Knaster--Tarski argument on a monotone map, standard machinery on a nonstandard carrier. What is claimed as bespoke, exactly: the carrier, the pruning-compatibility of the congruence, and the properness base case --- nothing else.

Three orientation remarks. \emph{Against the modal face}: the counterfactual companion's perturbation cycle --- reopen, restate, re-prune --- operates on exactly the objects constructed here; its reopening operator is a move in $T$, its re-pruning is passage to the reduced form, and its refusal to refine an ordering is, in this paper's vocabulary, the refusal to impose structure on $M$ beyond what the congruence requires. The complementarity is intended literally: one object, two instruments.

\emph{Against the semiring literature}: commutative semirings and their congruences are classical \cite{golan1999semirings}; what is framework-specific is the two-layer architecture --- graph structure intrinsic and foundational, weights extrinsic and customizable --- and the identity-plus-quotient mechanism through which the second layer inherits the first.

\textit{No additive inverses: the constraint, by layer.} The no-additive-inverses constraint is likewise not an aesthetic, and it splits by layer, settled at both. At the unweighted layer it is \emph{forced outright}: addition on $U$ is the informational union, hence idempotent ($x + x = x$), and an additively idempotent structure with inverses is trivial --- add $-x$ to both sides of $x + x = x$ and get $x = 0$, collapsing $U$ to a point. Two lines, airtight.

\textit{The weighted layer: three corrected claims.} At the weighted layer, three corrected claims replace a refuted one. (a) $T$ becomes a ring exactly when $M$ has additive inverses: the if direction is the display, $(-a)A + aA \equiv 0 \cdot A$ under the same-Aethus clause and the deletion convention above (§\ref{sec:extrinsic}); the only-if direction and its elementwise sharpening --- a nontrivial inverse pair in $T$ forces one in $M$, so zerosumfree $M$ leaves $T$ with no nontrivial inverses --- both immediate in the $M[U]$ coordinatewise reading (§\ref{sec:placement2}): $x + y \equiv 0_T$ means $x(A) + y(A) = 0$ coordinatewise, exhibiting a zero-sum pair in $M$ directly. (b) Zerosumfreeness is the guard \emph{simpliciter}: it is what makes the algebraic preorder an order and the weighted-parent relation a nontrivial one. The non-negativity rider in the weighted-parent definition is \emph{redundant} under the paper's own order stipulation --- $0 \leq \mu$ holds for every $\mu \in M$ via $\mu = 0 + \mu$, so every coefficient is non-negative definitionally, zerosumfree or not; the rider is annotated as automatic where it appears and retained for readability across general cones. The operative distinction remains zerosumfreeness versus the mere availability of a distinguished sub-cone ($\mathbb{R}$ contains $[0, \infty)$ yet is a ring). (c) The conditional forcing: once inverses are admitted, parenthood is total ($z = x + (-y)$ always exists, and the rider --- already vacuous per (b) --- cannot block it), $S(x)$ balloons, and the reduced-form order trivializes; zerosumfreeness (§\ref{sec:extrinsic}) is what the constraint protects, forbidding the inverses at the root.

One identity travels from the fuller development with its weighted scope characterized under a local hypothesis, and it reads as a corollary of this layer split: the \emph{cancellation identity} --- were subtraction admitted, $(X - Y)\cdot XY = X^2Y - XY^2 = XY - XY = 0$ for unweighted $X, Y$ by idempotence, the super-contradiction collapsing to the absorbing element by pure factoring. At the weighted layer the scope is precisely single carriers at equal scalars: $A = \alpha a$ and $B = \beta b$ give $A^2B = \alpha^2\beta\,(ab)$ and $B^2A = \alpha\beta^2\,(ab)$, agreeing iff $\alpha = \beta$ or a scalar vanishes --- under multiplicative cancellativity on nonzero elements, a \emph{local hypothesis} rather than a stipulation on $M$ (it holds at the base case and subsumes the zero-divisor hedge on the perturbation clause) --- while general combinations fail --- $x = \alpha(a_1 + a_2)$, $y = \alpha b$: the cross-carrier term $2\alpha^3\,a_1a_2b$ of $x^2y$ has no counterpart in $y^2x$ --- equal scalars restoring at single carriers the symmetry idempotence supplies at the unweighted layer.

% >>> ADDED (Fable 5, August 26, 2026): the complement/implication/truth cluster's algebra face,
% per the author's direction; the source cluster is the fuller development's, installed there
% August 19, 2026 at the subtraction principle.
\paragraph{The cluster the identity belongs to: complement, implication, and truth, as algebra.}\label{par:complement} The cancellation identity is one face of a three-part package recently consolidated in the fuller development \cite{Aethus}, and the package's natural home in \emph{this} paper's vocabulary is the quotient by validity equivalence, where each part becomes a one-line identity of $\bar{T} := T/\!\equiv_V$, writing $[x]$ for classes and $\bar{0} = [\Hyp]$ for the zero. The construction consumes one import, and it must be stated with care. The fuller development's partition identity refines a class over an attribute's alternatives, but in $\bar{T}$ it is not an equality: validity equivalence kills invalid carriers and identifies nothing else, so for the alternatives $\{b_j\}$ of an attribute stated in $B = b_1$'s terms, $[A]$ and $\sum_j [A\, b_j]$ have coefficients $1$ and $0$ at $A$. It enters instead as an \textbf{expansion}, $E_B[A] := \sum_j [A\, b_j]$, which is additive and multiplicative on $\bar{T}$: a cross term $[A\, b_j][C\, b_k]$ with $j \neq k$ dies by the rivals' joint invalidity, and $b_j b_j = b_j$ keeps the diagonal, so $E_B[A]\,E_B[C] = E_B[AC]$; and it respects the congruence, an invalid carrier's branches being invalid. The expansion is an identity, $E_B[A] = [A]$, exactly when $A$ is determinate on the attribute, the carving held; otherwise the fuller development's identity holds for the valuations that are additive across the attribute's alternatives, $v([A]) = v(E_B[A])$. The \textbf{expansion-relative complement of $B$ with respect to $A$}, written $R_B(A)$, is then the explicit residue
\begin{equation}
R_B(A) \;:=\; \sum_{j \geq 2} \,[A\, b_j] \qquad (\text{the name } A - AB \text{ is used for it below, and only as a name}),
\end{equation}
the branches of $A$ that $B$ kills, satisfying $E_B[A] = [AB] + R_B(A)$ and not in general $[A] = [AB] + R_B(A)$, since $E_B[A] \ne [A]$ in general (with $A = \mathbf{1}$ and two valid exclusive alternatives, $E_B[\mathbf{1}] = [b_1] + [b_2] \ne [\mathbf{1}]$), so that the subtraction is licensed from the expansion and not from the coefficientwise order on $[A]$ --- \emph{constructed}, not solved for: existence is by the expansion, no additive inverse is invoked anywhere, and the zerosumfreeness this section protects stands untouched, the difference notation licensed exactly as the order's residues are licensed above. (Pre-quotient, in $T$ itself, no single element deserves the name: the struck invalid summands remain standing as distinct tagged terms --- the tags-standing design of \S\ref{sec:extrinsic} --- and ``the'' complement exists only where the congruence has sent them all to $\bar{0}$, which is the algebraic content of the fuller development's clause that the complement is defined \emph{up to validity equivalence} and at no finer tier.) Two derivations then close the circle. The traveling identity is re-derived rather than re-imported: $(A - AB)\cdot[AB] = \sum_{j \geq 2}\,[A^{2}\, b_j\, b_1] = \bar{0}$, each term dying because rival alternatives of one attribute are jointly invalid, the carrier's invalidity annihilating the class whatever weight $A^{2}$ carries --- the complement and the residue disjoint by construction, which is what the cancellation identity was saying from the other side. (And note the hypothesis bookkeeping: the scalar-symmetry caveat above belongs to the \emph{pre-quotient} equality $X^2Y = XY^2$ in $T$; at the quotient, carrier invalidity does the annihilating alone, and the identity holds with no scalar hypothesis at all.) And \textbf{implication} arrives as the identity the fuller development takes as its definition, $A \Rightarrow B \,:\Longleftrightarrow\, [A][B] = [A]$, (multiplying a weighted parent in prunes nothing), whereupon the expansion makes the two-sided theorem immediate: \begin{equation} A \Rightarrow B \;\overset{(\mathrm{C})}{\Longleftrightarrow}\; A - AB = \bar{0}, \qquad\qquad A \Rightarrow \lnot B \;\overset{(\mathrm{C})}{\Longleftrightarrow}\; A - AB = E_B[A], \end{equation} the left because the killed branches carry the whole difference $E_B[A] - [A][B]$, so they all vanish exactly when $E_B[A] = [AB]$, the implication read at the expansion. $[A][B] = [A]$ gives that at once, the rivals of a stated $b_1$ being invalid; the converse needs a \emph{closure rule}, marked (C) in the display, which the tags do not supply by themselves: that an Aethus whose rival branches are all invalid states the surviving branch. Without (C) only the forward implications hold. The fuller development's deductive closure supplies that rule where the content generates the invalidity; for an externally tagged invalidity it is an added assumption, without which $[A]$ and $[A b_1]$ stay distinct classes of $\bar{T}$ while the residue vanishes, and the left equivalence holds from left to right only --- the difference notation, and the forward directions' canceling of $[AB]$, both licensed by $\bar{T}$'s additive cancellativity, inherited pointwise per the contracted-semiring identification of \S\ref{sec:placement2} --- the right because $[A][B] = \bar{0}$ leaves every branch killed. Its right-hand side is $[A]$ itself exactly when $A$ is determinate on the attribute, which is where the fuller development's partition identity holds as an equality. The classical equivalence $A \to B \iff A \wedge \lnot B \equiv \bot$ is thereby a theorem of the residue rather than a stipulated connective law, no negation primitive appearing: what would have been $\lnot B$ is carried by which joints the congruence annihilates. \textbf{Truth} is the remaining reading of the same product: $A$ is \emph{true relative to} $B$ iff $[A][B] \neq \bar{0}$. And the explosion principle, defused in the fuller development on the implication side, is here the two readings of one equation. $\bar{0}\cdot[x] = \bar{0}$ read against the implication identity says a contradiction implies everything --- $\bar{0} \leq [x]$ for every $x$ under the algebraic order, the hyperinvalid class everyone's child, the order's bottom --- while the same equation read against the truth test says the contradiction is true relative to nothing. Structurally real and truth-inert at once: the separation that makes explosion harmless is one equation admitting two readings, never confused.
% <<< END ADDED
Three layers complete its picture: exactly, since $x^2y = x\,(xy)$ and $y^2x = y\,(xy)$, the identity is \emph{agreement-in-action on the joint assertion}; the solution set obeys a scale-invariance law --- $(\lambda x)^2(\lambda y) = \lambda^3\,x^2y$, closure under common rescaling --- the formal face of the argument that equal scalars could not fail, their failure singling out an absolute weight scale against the relational weight principle \cite{Aethus}; and the exhaustive classification is conjectured and booked rather than claimed: the trivial family with \emph{either factor} $0_T$ (both sides vanish by absorption; the identity is swap-symmetric, so the families read accordingly), together with $x = y$ and the single-carrier equal-or-vanishing-scalar pairs, perturbations $y = x + w$ dying by zerosumfreeness ($x^2w + w^2x = 0$ forced, then killed termwise, for measure types without zero divisors), the incomparable case awaiting independent verification, a layered commitment the physics-face companion advertises as falsifiable.

The correspondence the physics-face companion needs is now specifiable: its refinement is declared descent as nature declares it, a blank a record separates holding the carving the record carries and splitting as an exclusive blank does, a record-free agreeing superposition irreducible, declaring no split until a carrier grants access, whereupon its splits are tagged, and so is it. The correspondence is intended over that whole declaration, the embedding of the companion's information states being the declaration nature's Aethus carries, read off record accessibility; that the two validity predicates agree on that image is the correspondence theorem still owed, and it is booked open here, this paper's normal form being the side it must be proved from; the physics-face companion records the same item from its side \cite{AethusQuant}. \emph{What remains open here}: the selection of the physical valuation is not this paper's question, and that valuation is not a congruence-respecting homomorphism, the Born weights not being one, as explained above, but the quantitative axiom of the physics face \cite{AethusQuant}, flagged and priced there; the categorical translation is booked as the forward note of §\ref{sec:extrinsic}, and the non-Archimedean stratification of $M$ is capacity in reserve.

\textit{The two inventories.}  Stipulated axioms, distinct from obligations: $M$ a commutative semiring with unit; on its additive side, zerosumfreeness (the guard, §\ref{sec:extrinsic}) and additive cancellativity, the algebraic order being a definition and the positive cone derived; on $T$'s quotient, the zero-term deletion convention; on the pair $(M, \ll)$, the relative-visibility axioms (the negligibility relation, strict dominance: transitivity; asymmetry; $0 \ll \nu$ for $\nu \neq 0$; relative downward closure; second-argument monotonicity; relative additive closure; scale-covariance for non-zero-divisor $\lambda$; composition (M-comp), \eqref{eq:mcomp}; sum-splitting where the two visibility tests must agree; and the base case recovering class $= \{0\}$ at $\mathbb{R}_{\geq 0}$), with the unified visibility rule and the $S(x)$ anchor, verification remaining open; and \emph{tag-extra-deductivity} on the consequence relation, jointness-disallowedness carried by tags, invisible to deductive closure (the clash-grade commitment, Definition~\ref{def:proper}). Obligations, distinct from the stipulations: coherence of component weights across admissible splitting expansions; the representation theorem of \S\ref{sec:child-proper-child}, a structure carrying the splitting identities, in which by Proposition~\ref{prop:splitting-obstruction}, under the retained idempotence, bilinearity and coefficient-uniqueness assumptions, a carrier can be identified with a splitting of exclusive cells only as their unit-coefficient sum, nonuniform weights remaining on distinct, non-idempotent weighted Aethae; the intrinsic operations on normalized support objects, union as the support union, with compatibility across decompositions; the correspondence theorem with the physics face; verification of the unified visibility rule. For the proper-child order $\le$, Theorem~\ref{thm:proper-lattice} supplies joins on the cone a fixed finite family of closure-independent carvings generates at real weights, that cone being a join-semilattice; meets exist only for compatible fillings of positive weight, and the bottom-adjoined completion is an order-theoretic device. The general case, carvings that are not closure-independent, acquired below a cell, or weighted beyond the real numbers, remains an obligation.



\begin{lemma}[{Pointwise order on $M[U]$}]\label{lem:pointwise}
The weighted-parent order on $T$, written $y \preceq_{\mathrm{add}} x$ to keep it visibly distinct from proper childhood $\le$ on carriers, is the \emph{pointwise} algebraic order on coefficients: $y \preceq_{\mathrm{add}} x$ iff $y(A) \leq x(A)$ for every $A \in U$ (coefficients total functions, zero off support). Forward, project the parenthood residue coordinatewise; backward, choose a residue at each coordinate over the finite support, which needs no cancellativity; cancellativity buys \emph{uniqueness} of each pointwise residue, hence canonicity of the assembled one. Corollary: $\preceq_{\mathrm{add}}$ on $T$ is a partial order iff $M$'s algebraic order is, so, under the standing additive-cancellativity assumption, zerosumfreeness is what makes the reduced-form order an order at both levels.
\end{lemma}

\begin{lemma}[The base case $M = \mathbb{R}_{\geq 0}$]\label{lem:basecase}
For real weights, the reduced form of $x$ is its valid-support truncation $x_V$ whenever $x_V \neq 0_T$: $x_V$ is proper; $x = x_V + x_I$ witnesses parenthood; and any nonzero proper $y \preceq_{\mathrm{add}} x$ vanishes on invalid carriers, hence $y \preceq_{\mathrm{add}} x_V$ pointwise (Lemma~\ref{lem:pointwise}). $S(x)$ is empty exactly when every positively weighted carrier is invalid  so \emph{maximum-existence discharges at the base case}. Further, evaluation at the blank coefficient, $x \mapsto x(\Null)$, lands in $M$, at the base case a homomorphism to $\mathbb{R}_{\geq 0}$ and in general via any $M \to \mathbb{R}_{\geq 0}$ homomorphism, and is multiplicative ($A \cdot B = \Null$ forces $\mathrm{Cl}(\mathrm{rep}\,A) \subseteq \mathrm{Cl}(\mathrm{rep}\,A \cup \mathrm{rep}\,B) = \mathrm{Cl}(\emptyset) \subseteq \mathrm{Cl}(\mathrm{rep}\,A)$, so $A = B = \Null$), respects the invalidity congruence by the kernel argument, a homomorphism's kernel being a congruence that contains the generating pairs ($1 \cdot X \mapsto 0$ for invalid $X \in U$, given the blank's validity), the truncation characterization being consumed instead by the stronger fact that evaluation also kills postulate-invalid \emph{combinations} (the scoped harmonization at Definition~\ref{def:reduced}) --- and separates when the blank is valid, as whenever tagging is non-pathological ($\Null$ untagged, base tags needing stated content), so \emph{properness likewise discharges at the base case}. The tempting total-valid-weight functional, by contrast, fails multiplicativity, valid times valid being taggable (the two bare slit-events).
\end{lemma}



\paragraph*{Provenance.} Extracted from the long companion's weighted-formulation section \cite{Aethus}, constructions and proofs intact (the name Aethus restored in September 2026), and revised under multi-model adversarial review directed by the author, whose rulings fixed the descent relations, the comparative proper-child condition, the zeros' layering, the doom--tag scope and the Born valuation's limits; the congruence note and fixed-point lemma were AI-drafted and author-approved, AI review served as pre-submission audit only, human refereeing precedes submission, and the audit trail is in the session record.


\appendix
\section{Countermodels to the superseded descent rule}\label{app:countermodels}

\noindent The cases below were computed under an earlier descent rule, on which an untagged Aethus could be doomed. They are kept because they show why the current rules are as they are; they describe that superseded rule, not the current model of \S\ref{sec:model}.

Two attributes $a$ and $b$, each with the values $1$ and $0$. The working universe holds the blank Aethus $\Null$, the four single-valued Aethae fixing one attribute, the four fixing both, the all-statements Aethus $\Hyp$, and, for the second and third cases, an Aethus $b^{\ast}$ holding $b$ in agreeing superposition. Every blank is exclusive except $b$ in $b^{\ast}$. Three cases, each checked exhaustively by the script that accompanies the paper.

\noindent\noindent\textit{The three cases that follow are the construction under the superseded descent rule, on which an untagged Aethus could be doomed; they are kept as the countermodel that forced the current rules, under which, as the paragraph on regimes as content states, $D$ is exactly the tagged carriers and the bare blank is valid. They are to be audited as that countermodel and not as the current model.} \emph{The ordinary blank.} Tag $a1b1$ and $a1b0$, so that $a1$ is untagged. Its declared proper children are $a1b1$, $a1b0$ and itself, and each has a tagged declared descendant, so $a1 \in D$: an untagged Aethus is doomed, because every way of filling its blank is impossible. Over the lattice order instead, $\Hyp$ lies beneath every Aethus, so $D$ would be all of $U$, the blank included; over declared descent the blank is valid, $a0b0$ being an untagged completion, and $D = \{\Hyp, a1b1, a1b0, a1\}$ together with the direct contradictions. At the weighted level $2 \cdot a1 + 3 \cdot a0b0$ has the reduced form $3 \cdot a0b0$, and $2 \cdot a1 + 5 \cdot a1b1$ has none.

\noindent\emph{The irreducible agreeing superposition.} $b^{\ast}$ declares no split of $b$: rule (c) gives it no declared proper child through $b$. Its declared descendants are itself and its fillings of $a$, none of which splits $b$, so nothing beneath it is tagged on $b$'s account, and $b^{\ast}$ is valid iff untagged. Tagging the fillings $b1$ and $b0$ changes nothing, since they are not beneath it.

\noindent\emph{The detector.} Multiply $b^{\ast}$ by a carrier $r$ recording $b$'s value. The product $b^{\ast} \cdot r$ carries access to its split, so by rule (c) it declares the children fixing $b$ at $1$ and at $0$, and each lies in $\mathrm{Base}$, a determinate $b$ contradicting the agreement. Every declared descendant of $b^{\ast} \cdot r$ that fixes nothing further has a tagged declared descendant, so $b^{\ast} \cdot r \in D$, and by the second postulate the system defaults to the disagreeing superposition $b1 + b0$, whose terms are valid: the double slit with a detector running. The same rule that dooms $a1$ dooms the recorded agreeing superposition, and leaves the unrecorded one alone.

\noindent\emph{The model, completed.} The three cases above are fragments of one fully specified model, computed exhaustively by the script that accompanies the paper. Its attributes are $X$ with pins $1, 2, 3$ and $b$ with pins $0, 1$, and a pin $e$ stating that $b$ is filled by exactly $\{0, 1\}$ and by no narrower filling. Its items are the states of proper nonempty fillings, the exclusions, the single-pin carvings, and $e$; its closure runs narrowing, the exclusions an exact filling carries, elimination, semantic mining of the single-pin carving, and invalidation. The relation structure declares $[X : \{1\}]$ jointly invalid with each determinate state of $b$, and $e$ jointly invalid with each, so that $b$ held with $e$ is the irreducible agreeing superposition $b^{\ast}$. The generators are the blank Aethus, the determinate and two-pin states of $X$, the determinate states of $b$, and $e$; the working universe is their closure under joint assertion and join. It holds 900 carriers, 752 of them tagged, invalid carriers remaining distinct by content. Carriers that share a filling can differ, the closure keeping every asserted state and deriving no wider one, so that $[X : \{1\}]$ and the joint assertion of $[X : \{1\}]$ with $[X : \{1, 2\}]$ are distinct Aethae. \paragraph{Regimes as content.} The construction above reads an unstated attribute as an exclusive blank by default, and that default is what opens the gap between $D$ and $I$: the doomed carrier $X_1$, with $b$ unstated and so read as exclusive, times $b^{\ast}$ is the valid double slit. The author's reading removes the default. A filling puts an attribute's open pins on the table and says nothing about validity below it; a blank declares a split only through what is intersected with it. Intersected with the single-pin carving $c_b$, the Aethus is the union of the two determinate states and splits into them; intersected with $b^{\ast}$, its determinate proper children are tagged; bare, it declares no split and is valid whenever it is not itself tagged. Rules (a)--(c) above are stated in this form, and Proposition~\ref{prop:doom-tag} follows: the third postulate's work on carved blanks is done by closure's own splitting rule, every doomed carrier is tagged, and $I = D$. On the completed model (\texttt{weight\_algebra\_model\_stated.py}), $|U| = 900$ with $782$ carriers tagged and $|D| = |I| = 782$; $X_1$ with $b$ unstated is valid, $X_1 \cdot c_b$ is tagged by the splitting rule, and $X_1 \cdot b^{\ast}$ is valid and outside $I$. An extension (\texttt{weight\_algebra\_broad\_check.py}) tests the case the completed model leaves untouched, a broad agreeing state whose determinate children are all excluded: it adds exclusions of a determinate $b$ by $[X : \{2\}]$, a marker $e_X$ holding $X$ irreducibly over $\{1, 2\}$, and a record of $X$'s value carrying $X$'s carving. With fillings passive, $|U| = 1100$ and $|D| = |I| = 1006$, every doomed carrier tagged; $[X : \{1, 2\}] \cdot c_b$ is valid, its product with $c_X$ is tagged, its product with $e_X$ is valid, and $e_X$ with the record is tagged. Splitting the broad agreeing state pin by pin instead reopens the gap: three doomed carriers go untagged, $[X : \{1, 2\}] \cdot c_b$ among them, and their products with $e_X$ are valid while lying in $I$. That is why a filling must declare nothing. The scripts named in this section (\texttt{weight\_algebra\_model.py}, \texttt{weight\_algebra\_model\_stated.py}, \texttt{weight\_algebra\_broad\_check.py}, \texttt{declared\_descent\_check.py}) accompany the paper as supplementary material; each runs with no arguments under Python~3 and prints the counts quoted here; they implement the closure under products and joins, the splitting rule and the passive-filling rule as defined here, and are to be re-run against any change to those definitions. One rule follows for the counterfactual semantics of the companion paper: reopening an attribute clears its filling to all pins and keeps its carving, so a reopened attribute is a scenario blank whose pockets are its determinate values, as the pocket construction assumes.

\textit{A countermodel of the superseded rule set, kept as such.} Read with the default this section first adopted, which the current rules (a)--(c) replace, declared descent splitting every blank left unstated and every broad agreeing state pin by pin, and giving $b^{\ast}$ no split, the model granting no access, the results are these: $D$ has 760 members, the 752 tagged carriers and eight untagged ones, every one of which fixes $X$ at $1$ and leaves $b$ open; $I$ has 768; and $I \setminus D$ consists of exactly eight carriers, every one of which fixes $X$ at $1$ and holds $b$ irreducibly, the carriers $X1 \cdot b^{\ast}$, doomed by nothing beneath them and killed by the quotient. With regimes as content, below, the gap closes. The joins behave as Proposition~\ref{prop:descent} states: $[X : \{1\}] + [X : \{2\}]$ holds the exclusion of $3$ and the carving $c_X$, and is not the blank; the join of all three determinate states holds $c_X$ and is not the blank; and $[X : \{1\}] + [X : \{2, 3\}]$ is the blank.

\printbibliography
\end{document}